\documentclass[11pt]{article}

\usepackage[margin=1in]{geometry}
\usepackage{amsmath}
\usepackage{amssymb}
\usepackage{hyperref}
\usepackage{booktabs}
\usepackage{listings}
\usepackage{xcolor}
\usepackage{parskip}
\usepackage{longtable}
\usepackage{array}

\lstdefinestyle{pystyle}{
  basicstyle=\ttfamily\small,
  keywordstyle=\color{blue!70!black},
  commentstyle=\color{gray},
  stringstyle=\color{orange!70!black},
  showstringspaces=false,
  breaklines=true,
  frame=single,
  language=Python,
}
\lstset{style=pystyle}

\title{DMRGPY Documentation}
\author{}
\date{}

\sloppy

\begin{document}
\maketitle
\tableofcontents

\section{Overview}

DMRGPY is a Python library for quasi-one-dimensional quantum lattice
models --- spin chains, spinless and spinful fermionic chains, Kondo
chains, parafermion chains, and bosonic chains --- solved with matrix
product state (MPS) methods, principally the density matrix
renormalization group (DMRG). Most calculations can also be run with
exact diagonalization (ED) on small systems, so DMRG results can be
cross-checked against an exact reference.

The library exposes a single, backend-agnostic Python API
(\texttt{src/dmrgpy/}). The same physics --- Hamiltonian, observables,
sites --- is written once, and DMRGPY dispatches the actual numerical
work to whichever solver is requested or available:

\begin{itemize}
  \item \textbf{DMRG}, run in-process through a compiled C++ extension
    (\texttt{itensor\_version=2} or \texttt{3}, built against a vendored
    copy of the \href{https://itensor.org/}{ITensor} library);
  \item \textbf{DMRG}, run through a live, in-process Julia session using
    ITensors.jl (\texttt{itensor\_version="julia\_live"});
  \item \textbf{DMRG}, run in a from-scratch, pure-Python/NumPy/SciPy
    reimplementation of the ITensor v3 API subset DMRGPY needs
    (\texttt{itensor\_version="python"}) --- no compiler or external
    language runtime required at all;
  \item \textbf{ED}, a pure-Python fallback for small systems
    (dense/sparse matrices via NumPy/SciPy), used automatically whenever
    a requested compiled backend is not available, and useful on its own
    as a correctness reference for DMRG.
\end{itemize}

\section{Installation}
\label{sec:install}

There are two independent install paths, and they cover different
backends:

\begin{itemize}
  \item \textbf{\texttt{pip install dmrgpy}} (via the repo's
    \texttt{pyproject.toml}) installs the pure-Python part of the
    package only --- the ED backend and the pure-Python
    \texttt{itensor\_version="python"} DMRG/TDVP backend both work
    immediately, no compiler needed. The compiled C++ backends
    (\texttt{itensor\_version=2}/\texttt{3}) are \emph{not} on PyPI:
    they're built against a $\sim$400MB vendored copy of ITensor, and
    \texttt{pyproject.toml} excludes \texttt{mpscpp2}/\texttt{mpscpp3}
    from the distributed package entirely. Without them, DMRGPY still
    works end to end --- it just falls back to ED/\texttt{"python"} (see
    mode.py) --- but if you want the compiled backends you need the
    git-clone path below instead.
  \item \textbf{A git clone + \texttt{python install.py}} is required
    for the compiled C++ backends (and is how this repository is
    developed). DMRGPY is used directly from \texttt{src/}, added to
    \texttt{PYTHONPATH} (or symlinked into \texttt{site-packages}) by
    the installer.
\end{itemize}

\begin{lstlisting}[language=bash]
python install.py                        # compiles ITensor v3 (mpscpp3, the default) + its pybind11 extension
python install.py --itensor-version=2     # compiles ITensor v2 (mpscpp2) instead
python install.py --itensor-version=both  # compiles both v2 and v3 backends
python install.py --gpp=g++-6             # use a specific compiler (needs g++ >= 6, LAPACK/BLAS)
python install.py --doctor                # check build requirements only, skip the build
python install_julia.py                   # alternative/additional: Julia backend
\end{lstlisting}

\texttt{install.py} first verifies every build requirement by actually
trial compiling and trial linking it (a C++ compiler, LAPACK/BLAS,
\texttt{pybind11}, \texttt{make}) rather than just checking version
strings, auto-detecting a working compiler (preferring a conda-provided
one when running under conda Python, since the extension shares a
process with conda's own numpy/scipy). Only once every requirement is
confirmed does it compile the vendored ITensor static library and the
pybind11 extension for the requested backend version(s). If nothing is
compiled at all (the extension build fails, or the user only ran
\texttt{install\_julia.py}), DMRGPY still works:
\texttt{itensor\_version="python"} requires no compiled extension, and
every call transparently falls back to ED when a requested compiled
backend isn't present.

No compiler or pybind11 is needed at all to use
\texttt{itensor\_version="python"} --- only NumPy/SciPy/SymPy/\texttt{numba}
(all four are core \texttt{pip install dmrgpy} dependencies;
\texttt{numba} is required rather than optional because
\texttt{algebra/kpmextrapolate.py} imports it at module level, and that
module is reached unconditionally from \texttt{manybodychain.py} via
\texttt{dynamics.py}/\texttt{kpmdmrg.py} --- this is separate from
\texttt{pyitensor}'s own opt-in JAX/numba-accelerated matvec kernel,
which \emph{is} genuinely optional and off by default, see
Section~\ref{sec:perf}). \texttt{jax} is optional either way.

\section{Quick start}

\begin{lstlisting}
from dmrgpy import spinchain

spins = ["S=1/2" for i in range(30)]      # 30-site spin-1/2 chain
sc = spinchain.Spin_Chain(spins)          # create the chain object

h = 0                                     # build the Hamiltonian symbolically
for i in range(len(spins) - 1):
    h = h + sc.Sx[i] * sc.Sx[i + 1]
    h = h + sc.Sy[i] * sc.Sy[i + 1]
    h = h + sc.Sz[i] * sc.Sz[i + 1]
sc.set_hamiltonian(h)

print("Ground state energy:", sc.gs_energy())
\end{lstlisting}

Every calculation follows this pattern: build a chain object for the
model of interest, build a Hamiltonian (and any observables) out of the
chain's own operators, call \texttt{set\_hamiltonian}, then call one of
\texttt{Many\_Body\_Chain}'s methods (\texttt{gs\_energy},
\texttt{get\_gs}, \texttt{vev}, \texttt{get\_excited},
\texttt{get\_dynamical\_correlator}, time-evolution methods, \ldots).
Most methods accept a \texttt{mode="DMRG"|"ED"} keyword so a result can
be cross-checked against exact diagonalization on small systems, e.g.:

\begin{lstlisting}
print("Energy with DMRG:", sc.gs_energy(mode="DMRG"))
print("Energy with ED:",   sc.gs_energy(mode="ED"))
\end{lstlisting}

The \texttt{examples/} directory contains 100+ self-contained scripts,
one per physical model or feature, organized into thematic subfolders
(\texttt{groundstate/}, \texttt{staticcorrelators/},
\texttt{dynamical\_correlator/}, \texttt{spin\_models/},
\texttt{fermion\_models/}, \texttt{boson\_models/},
\texttt{time\_evolution/}, \texttt{excited\_states/},
\texttt{entanglement/}, \texttt{topological/}, \texttt{kondo/},
\texttt{non\_hermitian/}, \texttt{finite\_temperature/},
\texttt{magnetization/}, \texttt{backend\_comparison/},
\texttt{utilities/}, \texttt{algebra/}, ...), and doubles as the
project's de facto regression suite (there is no separate unit-test
suite for the Python code). See \texttt{examples/readme\_examples/} for
the snippets shown in \texttt{README.md}, and
\texttt{examples/backend\_comparison/} plus the \texttt{v2\_VS\_v3\_*}-named
examples living inside each theme folder (e.g.
\texttt{examples/spin\_models/v2\_VS\_v3\_heisenberg\_spin\_half}) for
scripts that directly compare backends against each other on
correctness, and \texttt{examples/groundstate/backend\_timing\_gs\_energy/}
for timing.

\section{Architecture}

\subsection{Entry-point classes}

Each physical model is a thin subclass of \texttt{Many\_Body\_Chain}
(\texttt{manybodychain.py}), which holds all shared state --- bond
dimension, sweep count, cutoffs, mode, temp folder, etc. --- and
implements the generic operations: \texttt{set\_hamiltonian},
\texttt{gs\_energy}, \texttt{vev}, \texttt{get\_gs}, correlators, time
evolution, entanglement entropy, and more. Model modules just define the
local Hilbert space and convenience operators for that statistics:

\begin{center}
\footnotesize
\begin{tabular}{lll}
\toprule
Module & Class & Notes \\
\midrule
\texttt{spinchain.py} & \texttt{Spin\_Chain} & alias of \texttt{Many\_Body\_Chain} itself \\
\texttt{fermionchain.py} & \texttt{Fermionic\_Chain} & spinless fermions; C/Cdag/N, JW string F \\
\texttt{spinfermionchain.py} & \texttt{Spin\_Fermionic\_Chain} & spinful fermions \\
\texttt{bosonchain.py} & \texttt{Bosonic\_Chain} & bosonic sites \\
\texttt{parafermionchain.py} & \texttt{Parafermion\_Chain} & parafermion sites \\
\texttt{mixedchain.py} & \texttt{Mixed\_Spin\_\allowbreak Fermion\_Chain} & spin \& spinful-fermion mix \\
\bottomrule
\end{tabular}
\normalsize
\end{center}

All of the above build \texttt{self.sites}, the list of per-site type
codes passed to \texttt{Many\_Body\_Chain.\_\_init\_\_}, from a
\emph{uniform} repetition of a single type code (e.g.\ \texttt{[0]*n}
for \texttt{Fermionic\_Chain}). Nothing in \texttt{sites.py}/the C++ or
pyitensor \texttt{Chain} constructors actually requires this --- the
site-construction layer (\texttt{mpscppN/get\_sites.h}'s
\texttt{SpinX(std::vector<int> const\&)},
\texttt{pyitensor/sites/siteset.py::SiteX}) already accepts an
arbitrary, heterogeneous per-site type-code list.
\texttt{bosonchain.py::SpinBoson\_Chain} and
\texttt{mixedchain.py::Mixed\_Spin\_Fermion\_Chain} are the two model
classes that build such a heterogeneous list directly (spin+boson and
spin+spinful-fermion respectively), each following the same pattern:
build every operator that could be meaningful at \emph{any} site, then
mask the invalid combinations to the literal int 0 per site.
\texttt{Mixed\_Spin\_Fermion\_Chain} represents a spinful-fermion
location as a pair of physical spinless-fermion (type-code 0) sites,
exactly like \texttt{fermionchain.Spinful\_Fermionic\_Chain}, so it
reuses the existing C/Cdag/A/Adag/F operator vocabulary and
Jordan-Wigner code unchanged. It currently only supports
\texttt{itensor\_version} \texttt{3} and \texttt{"python"}: a
Jordan-Wigner string that has to cross a non-fermionic (spin) site
relies on \texttt{SiteSet::op()}'s ``unrecognized F request resolves to
identity'' fallback, present in ITensor v3
(\texttt{mpscpp3/ITensor/itensor/mps/siteset.h}) and in pyitensor's own
\texttt{SiteType.matrix()} (\texttt{pyitensor/sites/base.py}), but not
in ITensor v2 (\texttt{mpscpp2/ITensor/itensor/mps/siteset.h}), which
would hard-abort instead.

\texttt{fermionchain.Spinful\_Fermionic\_Chain\_Native} uses site-type
code 1 instead (\texttt{HubbardSite}/\texttt{ElectronSite}, ITensor
v3's own stock 4-state site,
\texttt{mpscpp3/ITensor/itensor/mps/sites/electron.h}) --- one
tensor-network site per physical site instead of
\texttt{Spinful\_Fermionic\_Chain}'s two interleaved type-0 sites.
\texttt{mpscppN/get\_sites.h}'s site-type dispatch already had a branch
for code 1 (\texttt{SpinX}'s \texttt{HubbardSite} case, both
constructors) before this class existed; only \texttt{mpscpp3} (the C++
side) and ED (\texttt{pyfermion.mbfermion.MBFermion.get\_operator},
extended to understand
\texttt{Cup}/\texttt{Cdn}/\texttt{Cdagup}/\texttt{Cdagdn}/\texttt{Nup}/
\texttt{Ndn}/\texttt{Ntot} at a physical site index) actually wire it up
today --- \texttt{mpscpp2} never registered a Hubbard/Electron site, and
neither pyitensor nor \texttt{mpsjulialive} implement site-type code 1.
See \texttt{fermionchain.py}'s class docstring for why this alternative
is not a general-purpose speedup over \texttt{Spinful\_Fermionic\_Chain}
despite the halved site count (a directly measured result, not a
theoretical claim).

\texttt{bosonchain.Bosonic\_Chain} uses a \emph{range} of type codes
rather than a single one: $100+\text{dim}$ per site, dim being the
requested local Hilbert space dimension
(\texttt{Bosonic\_Chain(n, maxnb=[\ldots])}, default $\text{dim}=4$,
i.e.\ code 104). This lets \texttt{mpscpp3/get\_sites.h}'s
\texttt{BosonFourSite} (\texttt{extra/bosonfour.h}, generalized from a
hardcoded 4-level site to an \texttt{Args("MaxOcc",\ldots)}-driven
arbitrary occupation cutoff) build a site of the actual requested
dimension instead of always the fixed 4-level one regardless of
\texttt{maxnb} --- previously a real, silent bug: the DMRG session only
ever saw type code 104 no matter what \texttt{maxnb} was, so DMRG and ED
results silently diverged for any non-default value.
\texttt{itensor\_version=3} and \texttt{itensor\_version="python"} both
understand the extended $102..199$ range today --- pyitensor's
\texttt{siteset.py::\_site\_class()} routes it to
\texttt{sites/boson.py::get\_boson\_site(dim)}, a small factory that
builds (and caches) a \texttt{SiteType} subclass per requested dimension
the same way \texttt{BosonFourSite} is generalized on the C++ side,
since pyitensor's \texttt{SiteType} machinery is otherwise entirely
class-level (no per-instance state) and has no other way to
parameterize a site's dimension. \texttt{mpscpp2} and
\texttt{mpsjulialive} still only handle the single code 104, so a
non-default \texttt{maxnb} should be run under
\texttt{itensor\_version=3} (the default) or \texttt{"python"} for
cross-backend agreement. ED (\texttt{pyboson/boson.py::BosonChain})
always builds the exact requested per-site dimension regardless of
backend.

\subsection{Infinite chains (\texttt{infinitechain.py}, \texttt{pyitensor/idmrg.py})}

\texttt{Infinite\_Many\_Body\_Chain}/\texttt{Infinite\_Spin\_Chain}
(\texttt{infinitechain.py}) are a deliberately \textbf{independent}
object, not a \texttt{Many\_Body\_Chain} subclass:
\texttt{Many\_Body\_Chain}'s whole design (mode dispatch between
DMRG/ED, a fixed \texttt{self.ns}-length site list,
KPM/dynamics/entanglement/excited-states machinery) assumes a
\emph{finite} chain throughout, and none of that has any meaning for a
translationally-invariant, infinite unit-cell chain.
\texttt{itensor\_version="python"} (default) or \texttt{3} are supported
(constructing anything else raises \texttt{NotImplementedError}) --
\texttt{2} and \texttt{"julia\_live"} have no iDMRG port. The public
surface is deliberately narrow: \texttt{set\_hamiltonian},
\texttt{gs\_energy} (energy \emph{per site}, the physically meaningful
quantity for an infinite system), static \texttt{vev}/\texttt{correlator},
and (via a finite-window reduction to the existing finite-chain KPM
stack, see below) \texttt{kpm\_finite} -- no excited
states, no entanglement/entropy in v1.

\textbf{Backend dispatch, unlike the finite-chain \texttt{mode.py}, is
explicit rather than automatic}: \texttt{gs\_energy} branches directly on
\texttt{self.itensor\_version} rather than falling back silently (there's
no ED oracle for an infinite system to fall back \emph{to}).
\texttt{itensor\_version="python"} runs \texttt{pyitensor/idmrg.py}'s own
growing algorithm (see below) and keeps the resulting
\texttt{IDMRGResult} in \texttt{self.\_result} for
\texttt{vev}/\texttt{correlator} to reuse. \texttt{itensor\_version=3}
instead calls the compiled \texttt{mpscpp3} backend's
\texttt{Chain::idmrg\_ground\_state} directly -- a line-by-line C++ port
of the same algorithm (built from the start with fresh Index objects on
every automaton tensor to sidestep the duplicate-Index bug class
described below, rather than needing the same retrofit) -- for the
energy density only; it has no static-correlator support, so
\texttt{self.\_result} is left \texttt{None} on that path and
\texttt{vev}/\texttt{correlator} raise \texttt{NotImplementedError}
regardless of whether \texttt{gs\_energy} itself already ran
successfully. Benchmarked (S=1/2 uniform Heisenberg chain,
\texttt{maxm=30}, \texttt{maxiter=60}, \texttt{etol=1e-9}): v3 is
$\sim$2.6--2.7x slower than \texttt{"python"} at both \texttt{n\_uc=1}
and \texttt{n\_uc=2}, with energy-density agreement between the two
backends to $\sim$1e-8. This is the \emph{gapless} (critical) case, the
hardest for any Krylov-based local solve since the correlation length
diverges and the local 2-site Arnoldi solve genuinely needs close to
its full \texttt{krylovdim} to converge at every macro-iteration
(confirmed by direct instrumentation) -- for a \emph{gapped} model the
same local solve typically converges in far fewer Krylov vectors, and
v3 can end up \emph{faster} than \texttt{"python"} (confirmed directly
on a dimerized \texttt{n\_uc=2} chain: v3 0.08s vs \texttt{"python"}
0.23s at \texttt{maxm=30}, \texttt{maxiter=60}). This asymmetry is
\texttt{arnoldi\_smallest\_real}'s own \texttt{early\_tol} early-exit at
work (see its own doc comment): every pre-existing caller of that
shared Arnoldi routine (\texttt{nhdmrg\_one\_sweep}) opts out
(\texttt{early\_tol<0}, the default) and is completely unaffected; only
\texttt{idmrg\_local\_solve}'s own call opts in, checking the same
residual-tolerance criterion the between-restart check already
trusted, just checked incrementally (every 3rd Krylov-vector addition,
from \texttt{m>=8} on) instead of only once a full
\texttt{krylovdim}-size subspace has been built -- mathematically the
same "stop once already converged" logic the pre-existing "happy
breakdown" branch already relied on (an exactly-zero residual), just
generalized to a tolerance. \texttt{TagSet} string-tag parsing
(\texttt{"Site"}/\texttt{"Link"}/\texttt{"a"}) was also confirmed by
profiling to be a real, multi-percent cost on its own (repeated on
every Krylov iteration and every micro-step) and is now cached in
function-local statics instead of re-parsed per call.
\texttt{Chain::idmrg\_ground\_state} is reached by constructing a
\texttt{Chain} directly from the unit cell's own
\texttt{site\_types} (\texttt{cppext.get\_backend(3).Chain(self.site\_types)}),
the same low-level pattern \texttt{nhdmrg.py} uses for its own
compiled-backend calls (\texttt{H.to\_terms()} fed straight into
\texttt{self.\_session.<method>(...)}) -- terms are shifted from
\texttt{infinitechain.py}'s own 0-based \texttt{h\_intra}/\texttt{h\_inter}
site convention to the codebase's usual 1-based one via the same
\texttt{MultiOperator.to\_terms()} every other compiled backend call
already uses (no separate Jordan-Wigner transform, since iDMRG doesn't
support fermionic terms yet on either backend).

The Hamiltonian is built from
\texttt{SxL[i]}/\texttt{SxC[i]}/\texttt{SxR[i]}-suffixed operators
(previous/central/next unit cell, \texttt{i=0..n\_uc-1}) rather than
absolute site indices; \texttt{\_canonicalize\_hamiltonian} validates
every term touches at most two adjacent cells and shifts away the
\texttt{L} suffix (\texttt{L}$\to$\texttt{C}, \texttt{C}$\to$\texttt{R})
before storing, so the stored Hamiltonian is uniformly "intra-cell
\texttt{C} terms" + "\texttt{C}-to-\texttt{R} inter-cell terms", each
physical bond attributed to exactly one cell.

\texttt{pyitensor/idmrg.py} implements the actual growing/infinite-size
algorithm (White 1992, generalized to an \texttt{n\_uc}-site unit cell
per McCulloch 2008): \texttt{\_build\_automaton} builds the periodic
per-sublattice MPO tensor directly as a hand-rolled finite-state
automaton (not extracted from a compressed finite reference system -- an
earlier design that was tried and abandoned, see the module's own
docstring for the two independent bugs that approach hit), then
\texttt{idmrg\_ground\_state} grows two environments \texttt{HL}/\texttt{HR}
from \texttt{HL=HR=None}, one unit cell at a time, reusing
\texttt{dmrg.py}'s own \texttt{\_lanczos\_ground\_state} and
\texttt{kernels.make\_matvec} for each micro-step's local 2-site solve
(no separate Lanczos implementation). Static correlators are
reconstructed after convergence from the last macro-iteration's own
per-sublattice left-canonical tensors via the standard infinite-MPS
transfer-matrix formalism (dominant left/right fixed points of the
per-unit-cell transfer operator).

\textbf{A confirmed, fixed reliability bug affecting static correlators}:
each micro-step's local 2-site Lanczos solve used to always start from a
fresh random vector, on the (mistaken) reasoning that every micro-step
inserts brand-new physical sites so there is no previous local tensor to
warm-start from. That's true of the \emph{sites}, but once bond dimension
saturates at \texttt{maxm} the \emph{shape} of the local eigenproblem at a
given unit-cell position repeats macro-iteration to macro-iteration -- so
for a gapless/SU(2)-symmetric model (routinely a (near-)degenerate local
ground manifold), a fresh random restart let Lanczos land on an arbitrary
member of that manifold every macro-iteration: the reported \emph{energy}
still converged fine (variational, degeneracy-robust) but \texttt{U\_list}
kept jumping between different members instead of settling into one
self-consistent unit cell, corrupting every downstream static correlator
(confirmed directly via the model-agnostic
$\langle H_{uc}\rangle$-equals-\texttt{n\_uc*density} self-consistency
identity: 0.4--0.6 gaps for a uniform \texttt{n\_uc=1} chain, not shrinking
with \texttt{maxiter}, and not fixed by best-of-6 independent-seed
reruns). Fixed by threading each macro-iteration's own local ground vector
back in as the \emph{next} macro-iteration's Lanczos warm start
(\texttt{\_local\_two\_site\_solve}'s \texttt{x0\_warm} parameter) -- this
closes the gap by roughly an order of magnitude for \texttt{n\_uc=2}
(0.34--0.40 down to $\le$0.05 across 7 independent trials, several
essentially exact) but, confirmed directly, does \textbf{not} fix
\texttt{n\_uc=1} (still 0.4--1.8, \texttt{IDMRGResult.state\_overlap}
plateauing around 0.5--0.65 rather than approaching 1 even after 80
macro-iterations) -- isolated to \texttt{n\_uc=1}'s own structurally
unique micro-step, where \texttt{p\_L==p\_R} always (every
macro-iteration's two active sites are literally the same sublattice, via
\texttt{\_fresh\_physical\_copy}), unlike \texttt{n\_uc=2} where this
never happens. \texttt{IDMRGResult.state\_overlap} (the normalized overlap
between consecutive macro-iterations' local ground vectors,
\texttt{None} if unavailable) is a new, cheap diagnostic exposing this
directly, independent of the $\langle H_{uc}\rangle$ self-consistency
check. See \texttt{docs/user\_guide.tex}'s iDMRG section for the
user-facing version of this caveat.

\textbf{Currently limited to \texttt{n\_uc<=2}} (enforced at
\texttt{Infinite\_Many\_Body\_Chain} construction) -- the micro-step
loop's sublattice pairing (\texttt{p\_L=mstep},
\texttt{p\_R=n\_uc-1-mstep}) only produces two genuinely \emph{adjacent}
active sites for \texttt{n\_uc} in $\{1, 2\}$; for \texttt{n\_uc>=3} an
intermediate micro-step's two active sites are several real,
not-yet-inserted sites apart, and \texttt{\_build\_automaton}'s own
per-sublattice pending-channel content differs between them (not just
its count), so there is no way to even wire them together correctly
without a genuine redesign of the growth scheme -- flagged as a known
follow-up rather than attempted.

\textbf{A confirmed, fixed bug worth knowing about if this module is
touched again}: \texttt{idmrg\_ground\_state} re-mints a fresh mpo-axis
Index for \texttt{HL}'s and \texttt{HR}'s own environment tensor at the
end of \emph{every} micro-step, rather than reusing whatever Index
\texttt{\_build\_automaton} happened to label that channel space with.
This is required because \texttt{\_build\_automaton}'s
\texttt{boundary\_idx} pool has only \texttt{n\_uc} distinct Index
objects total, reused verbatim every time a given sublattice position
comes up again (every macro-iteration past the first) -- left
un-broken, \texttt{HL}'s and \texttt{HR}'s own mpo axes can end up being
the \emph{same} Index object (guaranteed for \texttt{n\_uc=1}, where a
single sublattice's left and right automaton legs are literally
identical), corrupting the effective Hamiltonian from the second
macro-iteration onward once both environments are non-trivial and used
together in the same local 2-site solve. Confirmed directly by
extracting the dense effective Hamiltonian at macro-iteration 2 of a
uniform \texttt{n\_uc=1} Heisenberg chain and finding its full
eigenvalue spectrum did not match exact 6-site ED at all, while
macro-iteration 1's spectrum matched to machine precision -- the same
self-referencing-Index hazard already seen once in this module
(\texttt{\_project\_channel}'s own docstring), here manifesting across
the growth loop's iterations rather than within a single tensor.
\texttt{idmrg.py}'s module docstring and the inline comments in
\texttt{idmrg\_ground\_state}/\texttt{\_relabel\_pos} have the full
derivation.

A separate, milder issue in the same neighborhood:
\texttt{U\_list[0]}'s own left bond and \texttt{U\_list[n\_uc-1]}'s own
right bond are, in a truly periodic MPS, the \emph{same} wraparound
bond, but each is independently truncated by its own micro-step's SVD --
nothing guarantees they come out numerically equal, and two independent
SVDs occasionally round a borderline singular value across the
\texttt{cutoff} threshold differently for what is conceptually one bond.
\texttt{\_dominant\_right\_fixed\_point} (used by both
\texttt{onsite\_expectation} and \texttt{two\_point\_correlator}) now
raises a clear \texttt{RuntimeError} when this happens instead of a
cryptic \texttt{reshape} crash; a smaller \texttt{maxm} that actually
binds (forcing both bonds to the same truncated dimension) avoids it,
see \texttt{tests/test\_infinite\_chain.py}'s own comment on this.

\textbf{Applying an MPO to the converged iMPS}: \texttt{idmrg.py} also
implements \texttt{apply\_mpo(result, W\_bulk, cutoff, maxdim)} -- the
infinite-chain analogue of \texttt{mpsalgebra.applyMPO}, and the
primitive a future iTEBD-style real/imaginary-time evolution feature
would build on. \texttt{grow\_by\_mpo} Kronecker-merges \texttt{W\_bulk}'s
own per-site bonds with \texttt{U\_list}'s (mirroring
\texttt{mpsalgebra.\_apply\_chain}'s per-site body, generalized to a
periodic index range including the wraparound cut), giving raw,
non-canonical tensors at bond dimension \texttt{chi\_A*chi\_W};
\texttt{\_canonicalize\_periodic} then re-canonicalizes/truncates them
back down via the standard two-sided fixed-point infinite-MPS procedure
(Or\'us \& Vidal, PRB 78, 155117 (2008)): factor the dominant left/right
transfer-matrix fixed points at every cut into square-root factors
\texttt{X\_p}/\texttt{Y\_p} (\texttt{\_psd\_sqrt\_factor}), SVD
\texttt{M\_p=X\_p@Y\_p} and truncate (reusing \texttt{svd.py}'s own
\texttt{\_truncate}), then build an \emph{asymmetric} gauge
(\texttt{G\_left[p]=U\_p\^{}H X\_p}, no \texttt{S} factor at all;
\texttt{G\_right[p]=Y\_p V\_p\^{}H S\_p\^{}\{-1\}}, the \emph{full}
inverse) -- asymmetric, rather than a textbook symmetric
\texttt{S\^{}\{-1/2\}} split on both sides, specifically so the result
comes out plain left-canonical (matching \texttt{IDMRGResult.U\_list}'s
own convention) with no separate bond-weight layer to thread through,
confirmed numerically before committing to the formula (a first,
symmetric-split derivation looked plausible by analogy to Vidal's
canonical form but reproduced Identity only to $\sim$1, not a small
residual). \texttt{\_all\_left\_fixed\_points} (the new, symmetric
counterpart to the existing \texttt{\_all\_right\_fixed\_points}) needed
one more fix beyond the obvious transpose-the-eigenproblem mirroring:
its own per-cut renormalization (needed to avoid blow-up across repeated
propagation, exactly like the existing right-fixed-point code) discards
the \emph{relative} scale between different cuts that
\texttt{\_canonicalize\_periodic}'s left-gauge construction turns out to
depend on (unlike the right-gauge one, which is provably
scale-invariant) -- tracked via an explicit accumulated-scale return
value and undone before use. A second, independent bug (an index-order/
transpose mismatch between \texttt{\_apply\_transfer\_from\_left}'s own
(ket,bra) convention and the (bra,ket) convention the isometry
derivation needs) was only found by deriving the exact identity the
gauge must satisfy by hand and checking it numerically term-by-term --
a genuine case where tests alone (which caught that \emph{something}
was wrong, via a hand-crafted internal left-canonicality check) weren't
enough to \emph{localize} the bug, only real algebra was. \textbf{Scope
restriction}: \texttt{W\_bulk} must be a \emph{bounded} (non-extensive)
periodic operator -- the Hamiltonian's own
\texttt{\_build\_periodic\_mpo}-built automaton is deliberately the
\emph{other} kind (an unconditional accumulator self-loop that needs a
genuine chain boundary to mean ``sum'', not ``product'') and is not a
valid \texttt{apply\_mpo} input; see \texttt{idmrg.py}'s own ``Applying
a (bounded) MPO to the converged iMPS'' section docstring for the
specific failure mode this would hit.

\textbf{Overlap between two converged iMPS}: \texttt{idmrg.py} also
implements \texttt{imps\_overlap(result\_a, result\_b, normalize=True)}
-- the thermodynamic-limit analogue of a finite MPS inner product. Since
a literal $\langle\phi|\psi\rangle$ over an infinite chain scales as
$\eta^N$ over $N$ unit cells (not a finite number unless $|\eta|=1$
exactly), the default return value is instead the \emph{per-site
fidelity} \texttt{eta\_ab/sqrt(eta\_aa*eta\_bb)}, where \texttt{eta\_ab}
is the dominant eigenvalue of the \emph{mixed} transfer matrix built
from the two states' own tensors (\texttt{\_transfer\_matrices}
generalized with a new optional \texttt{bra\_list} parameter, defaulting
to the self-overlap case every existing caller still uses) and
\texttt{eta\_aa}/\texttt{eta\_bb} are each state's own self-overlap (via
the existing \texttt{\_dominant\_right\_fixed\_point}, reused verbatim
-- valid since a self-overlap transfer tensor is always square). The
mixed case needed a dedicated eigen-solve,
\texttt{\_dominant\_eigenvalue\_mixed}, rather than reusing
\texttt{\_dominant\_right\_fixed\_point} directly: two independently
converged/truncated iMPS need not share a bond dimension, so the
composed mixed transfer tensor is
\texttt{(chi\_a,chi\_b,chi\_a,chi\_b)}, square only in the
\texttt{chi\_a==chi\_b} self-overlap special case
\texttt{\_dominant\_right\_fixed\_point} itself assumes.

\textbf{A second, real bug found while implementing \texttt{imps\_overlap}}
(needed a Hamiltonian with an onsite/field term to exercise a
meaningfully non-trivial ``orthogonal states'' test case, which
surfaced two bugs in \texttt{\_build\_periodic\_mpo}'s onsite-term
handling, both pre-dating and unrelated to \texttt{imps\_overlap}
itself): \texttt{\_onsite\_matrix}'s own unpacking loop expected
3-tuples (\texttt{rel\_site, coef, mat}) but \texttt{\_build\_periodic\_mpo}
had already consumed \texttt{rel\_site} as its \texttt{onsite\_by\_p}
dict key, storing only \texttt{(coef, mat)} pairs -- a
\texttt{ValueError: not enough values to unpack} for \emph{any}
Hamiltonian with an onsite term (dead code before this, since every
existing test used bond-only Hamiltonians). Fixing just that crash
would have been unsafe: the underlying automaton wiring was also wrong.
Onsite content was added on the accumulator's own \texttt{F,F}
self-loop (\texttt{Id + onsite\_mat}) rather than a direct
\texttt{S,F} transition -- no existing branch ever populates
\texttt{S,F} at all, so (1) for a bond-less Hamiltonian, \texttt{W}
stays block-diagonal between $\{S\}$ and $\{F,\ldots\}$ forever, and the
boundary-extracted $\langle S|W^N|F\rangle$ (exactly what
\texttt{\_project\_channel}'s boundary tensors compute) is identically
0 for every $N$ -- confirmed directly: a pure-field \texttt{n\_uc=1}
chain converged to energy density 0 and $\langle S_z\rangle\sim0$
regardless of field strength, instead of the exact, closed-form
$-|B|/2$/$\mathrm{sign}(B)\cdot0.5$; and (2) once a bond term \emph{did}
activate \texttt{F}, every further site absorbed while \texttt{F} was
already hot re-added the onsite content on top of an already-summed
running total (an exponential, multiplicative blow-up, not the intended
linear sum) -- confirmed directly: energy density reached $-10^{23}$ by
$B=0.3$ and $-10^{69}$ by $B=1.0$ on an otherwise-normal dimerized
\texttt{n\_uc=2} chain, with \texttt{.converged} staying \texttt{False}
throughout instead of a modest, field-dependent shift. The fix moves
the onsite content onto the direct \texttt{S,F} transition (leaving
\texttt{F,F} as plain \texttt{Id}) -- the standard textbook
automaton-MPO construction for summing an onsite term into a running
total (the same upper-triangular-in-channel-index structure any
finite-range Hamiltonian MPO derivation uses), verified both by a
by-hand matrix-product induction argument and numerically (the same
field tests above now reproduce the exact closed-form field-only result
and a finite, converging energy density with a bond term present).

\textbf{Direct sum of two converged iMPS}: \texttt{idmrg.py} also
implements \texttt{imps\_sum(result\_a, result\_b, cutoff, maxdim)} --
the periodic-chain analogue of \texttt{mpsalgebra.sum}.
\texttt{\_periodic\_direct\_sum} block-diagonally concatenates
\texttt{result\_a}'s and \texttt{result\_b}'s own bond spaces at
\emph{every} unit-cell cut, including the wraparound (unlike
\texttt{mpsalgebra.sum}'s finite chain, which has two dimension-1
boundaries to instead concatenate along -- a periodic chain has no such
boundary, so the sum is block-diagonal everywhere); the raw result is
then re-canonicalized/truncated by reusing \texttt{apply\_mpo}'s own
\texttt{\_canonicalize\_periodic} verbatim, no new gauge-fixing algorithm
needed.

\textbf{A real, load-bearing correctness issue found and fixed while
implementing this} (not merely a theoretical concern -- reproduced
directly before the fix): \texttt{\_canonicalize\_periodic}'s two-sided
fixed-point procedure assumes the combined transfer matrix has a
\emph{unique} dominant eigenvalue, picked via
\texttt{np.argmax(np.abs(w))} in \texttt{\_dominant\_right\_fixed\_point}.
Every \texttt{IDMRGResult} is individually normalized to self-overlap
eigenvalue \texttt{eta=1} exactly (left-canonical SVD construction), so
summing two \emph{different}, ordinarily-converged \texttt{IDMRGResult}s
-- the natural use case -- produces a combined transfer matrix with an
\emph{exactly degenerate} leading eigenvalue (two blocks, each eigenvalue
1). Before a guard was added, \texttt{np.argmax} silently picked
\emph{one} of the two tied eigenvectors (confirmed directly: summing two
independently-converged, oppositely Sz-polarized \texttt{n\_uc=1} product
states reliably collapsed to just one branch, with \emph{which} branch
survived sensitive to floating-point tie-breaking noise inside
\texttt{np.linalg.eig} -- not even reproducible run to run), silently
discarding the other branch entirely while still returning a
plausible-looking, correctly-normalized \texttt{PeriodicMPS}. This is the
classic ``looks right, isn't'' failure mode this codebase's own testing
culture exists to catch. The fix: a shared helper,
\texttt{\_check\_dominant\_eigenvalue\_nondegenerate}, checks the top two
eigenvalue magnitudes of the composed transfer matrix and raises
\texttt{RuntimeError} when their ratio exceeds $1-10^{-9}$. This
threshold is set far tighter than the loosest value that would have
passed every legitimate case tried -- the genuine tie this check exists
to catch lands at a $\sim 10^{-15}$ relative gap (exact to machine
precision), so $10^{-9}$ still rejects it with 6 orders of magnitude to
spare, while leaving a wide safety margin below: a gapped dimerized
\texttt{n\_uc=2} chain's top two magnitudes were $(1.0, 0.093)$; a
gapless/critical uniform \texttt{n\_uc=1} chain showed a clear
\texttt{maxm}-dependent narrowing (top-two gap $0.376$ at
\texttt{maxm=40}, $0.116$ at \texttt{maxm=60}, $0.160$ at
\texttt{maxm=80} -- not perfectly monotonic, since \texttt{n\_uc=1} has
its own known convergence limitation, see above, but never close to
tight) -- a threshold merely ``loose enough for today's test suite''
would risk a \emph{future} false positive on an ordinary, non-degenerate
correlator call at a larger \texttt{maxm} a user might reasonably pick
for better accuracy near criticality, so $10^{-9}$ was chosen
deliberately over a looser value like the originally-tried $10^{-6}$.
Physically, the tied case is a ``cat state'' superposition of two
macroscopically distinct branches; representing it as a single
injective/canonical periodic MPS is out of scope for this module's
existing (single-fixed-point) correlator machinery -- \texttt{imps\_sum}
therefore raises clearly there rather than silently returning one
arbitrary branch, mirroring \texttt{apply\_mpo}'s own ``bounded operators
only'' scope restriction in spirit. The shared helper is used by
\texttt{\_dominant\_right\_fixed\_point}, \texttt{\_dominant\_left\_fixed\_point},
and \texttt{\_dominant\_eigenvalue\_mixed} (the mixed-transfer-matrix
eigenvalue \texttt{imps\_overlap}'s own cross term uses) alike -- all
three previously used an unguarded \texttt{np.argmax}, the identical
latent bug, just not yet exercised by an existing caller for the
left/mixed cases -- so this same guard now protects
\texttt{onsite\_expectation}/\texttt{two\_point\_correlator}/
\texttt{imps\_overlap} for free too; confirmed not to trip on any
pre-existing test or on \texttt{imps\_overlap}'s own realistic (orthogonal
and gauge-comparison) cases (all previously-passing
\texttt{tests/test\_infinite\_chain.py} cases still pass).

\textbf{Excited states: the tangent-space/quasiparticle excitation
ansatz (\texttt{pyitensor/idmrg\_excitations.py}, a new module)}: an
infinite chain's excitations form a momentum-resolved band $E(k)$, not a
single discrete state a finite chain has, so
\texttt{Infinite\_Many\_Body\_Chain.excitation\_energies(k, n=1)}/
\texttt{excitation\_gap(ks=None)} implement the standard single-mode/
quasiparticle ansatz (Haegeman et al.) instead of anything reusing the
growing algorithm's own truncated environments: a tangent-space vector
$|\Phi_k(X)\rangle = \sum_n e^{ikn} |\ldots A\,A\,B_n\,A\,A\ldots\rangle$
with one excitation tensor $B(X)$ inserted at every unit-cell position,
weighted by momentum $k$. Multi-site unit cells are handled by first
grouping the \texttt{n\_uc} site/automaton tensors into one effective
supersite (\texttt{\_group\_unit\_cell}), so the rest of the module only
ever deals with a uniform, single-supersite chain; every 2-site bond
term must have reach exactly 1 after grouping (checked directly on the
grouped automaton's own channel structure, \texttt{\_check\_reach\_one},
rather than by re-inspecting the original Hamiltonian term lists). The
construction: a null-space gauge-fixing tensor $V_L$
($B(X) = \mathrm{reshape}(V_L X)$, $V_L^\dagger A = 0$, ensuring $B$ is
orthogonal to the ground state itself); $l=I$ exactly (left-canonical
$A$) and $r$ = the ordinary dominant right fixed point already computed
elsewhere in \texttt{idmrg.py}; two regularized (``energy density
subtracted off'') environments $L_h$/$R_h$ solved as dense linear
systems; and a momentum-dependent effective Hamiltonian $H_{\rm eff}(k)$
built from a \emph{finite} list of diagrams -- no genuine
infinite/geometric momentum sum is needed, since the gauge condition
makes any diagram where the ket differs from the bra by more than one
adjacent unit cell vanish identically (confirmed directly, not just
assumed, on a finite-ring tensor-network contraction), so only unit-cell
separations $0, \pm1$ ever contribute for this module's reach-1
Hamiltonians.

\textbf{Real, load-bearing bugs found and fixed while implementing
this} (via extensive cross-checking against a from-scratch finite-ring
tensor-network contraction, bypassing this module's own machinery
entirely -- not just internal self-consistency checks, which turned out
to be insufficient to catch these): (1) the energy density used to
regularize $L_h$ (\texttt{Source\_L}'s own trace) needs an
$r$-weighted trace ($\mathrm{tr}(r\, \mathrm{Source\_L})$), not the
plain trace \texttt{Source\_R}'s own regularization correctly uses --
confirmed directly by comparing a single-site
\texttt{apply\_transfer\_from\_left} step against
\texttt{idmrg.onsite\_expectation}'s already-validated formula, off by a
measurable, non-numerical-noise amount with the plain trace. (2) Three
of the effective Hamiltonian's seven diagrams (the ones ending in a
\texttt{\_cap\_left} step, threading an environment into the excitation
tensor from the left) were each missing a further
\texttt{\_cap\_right(..., r)} closure on the \emph{other} side -- easy
to miss because omitting the analogous closure is harmless whenever it
would be capped by $l$ (=Identity, a no-op), which is exactly the case
for the other four diagrams, but not harmless for $r$ (not the identity
in general) -- found by comparing individual diagrams against the
finite-ring reference and finding an exact, large (not roundoff-scale)
mismatch, fixed, and re-verified to match to $\sim10^{-13}$. (3) The
tangent-space norm's own metric (the fixed point $r$) is, in general, a
severely ill-conditioned matrix -- its eigenvalues are exactly the
ground state's own entanglement spectrum, and any weakly-entangled
(e.g.\ deep-paramagnetic) ground state has a steeply decaying one
(confirmed directly: eigenvalues $[0, 0, 0, 0.0073, 0.9927]$, condition
number $\sim10^{10}$, for a transverse-field-Ising ground state at bond
dimension 5) -- solving the naive generalized eigenproblem
$H_{\rm eff}(k)[X] = E\,(X r)$ directly against this metric gave
numerically garbage results; the fix (\texttt{\_reduce\_metric})
substitutes $X = \tilde X\, \mathrm{diag}(1/\sqrt{\mathrm{eig}})\,
V^\dagger$ restricted to $r$'s own well-conditioned eigenspace (the
dropped directions have exactly zero physical norm, not merely small,
so nothing physical is discarded), turning the problem into an
ordinary, well-conditioned Hermitian eigenproblem.

\textbf{A known, currently unresolved limitation, not silently papered
over}: after all three fixes above, the ansatz matches the exact
free-fermion single-magnon dispersion of a field-polarized XX chain
(bond dimension $D=1$, an uncorrelated/product-state ground state) to
$\sim$14 digits across the entire Brillouin zone -- but for a genuinely
entangled ground state ($D>1$), the resulting dispersion comes out
anomalously flat compared to the expected answer, despite every
individual diagram still independently matching the same finite-ring
reference to $\sim10^{-13}$ and $H_{\rm eff}(k)$ remaining exactly
Hermitian at every momentum tried. This was not root-caused despite
extensive investigation (ruled out: metric conditioning, per (3) above;
$L_h$/$R_h$'s own internal consistency, verified both via the linear
solve's own residual and via independent direct iteration of their
defining recursion; iDMRG's own convergence quality, which did not help
even at \texttt{state\_overlap}$\sim0.9999$).
\texttt{build\_excitation\_environment} therefore explicitly rejects
$D>1$ with \texttt{NotImplementedError} rather than silently returning a
dispersion known, in at least some cases, to be wrong -- see
\texttt{pyitensor/idmrg\_excitations.py}'s own module docstring (``KNOWN
LIMITATION'' section) for the full account, and
\texttt{examples/idmrg/excitation\_gap\_xx/main.py} for the validated
$D=1$ case.

\textbf{A cheaper, cruder alternative: the local superblock gap
(\texttt{pyitensor/idmrg.py::local\_excitation\_gap})}: rather than a
fresh tangent-space construction, this reuses the growing algorithm's own
final micro-step exactly as it already runs -- \texttt{idmrg\_ground\_state}
now stashes the last micro-step's local-solve ingredients
(\texttt{HL}/\texttt{HR} environments, the two MPO tensors, the two
physical Indices, the converged local ground vector/energy) on
\texttt{IDMRGResult.local\_superblock} -- and re-diagonalizes that
\emph{same} 2-site effective Hamiltonian (rebuilding its matvec via
\texttt{kernels.make\_matvec} on the stored pieces) for its second-lowest
eigenvalue instead of only its ground state, returning the difference.
Framed explicitly as the infinite-chain analogue of finite DMRG's own
Lagrange-multiplier/overlap-penalty excited-state method
(\texttt{dmrg.py}'s \texttt{dmrg\_excited}/\texttt{Chain.excited\_states}):
that method needs a penalty term threaded through a whole re-sweep because
it enforces orthogonality against a \emph{separate}, externally-held
ground-state MPS one local update at a time; here the ``state to stay
orthogonal to'' is just the local ground vector already found, in the very
same local Hilbert space, by the very same solve, so the constraint is
enforced \emph{exactly} via deflation ($P = I - |\psi_0\rangle\langle\psi_0|$;
$\text{deflated\_matvec}(v) = P(\text{matvec}(P(v)))$, Hermitian since both
$P$ and the underlying local Hamiltonian are) rather than approximately via
a finite penalty weight -- mathematically, a Hermitian operator's
constrained stationary points (extremize $\langle\psi|H|\psi\rangle$
subject to $\langle\psi|\psi\rangle=1$, $\langle\psi|\psi_0\rangle=0$) are
exactly its \emph{other} eigenvectors, so this is what a penalty method
converges to anyway as its weight $\to\infty$, obtained directly with
nothing to tune. Implementation mirrors \texttt{\_local\_two\_site\_solve}'s
own small-dimension fallback (dense \texttt{eigh} for $\dim\le3$) and
otherwise reuses \texttt{dmrg.\_lanczos\_ground\_state} (imported already)
on the deflated matvec, rather than introducing a second Lanczos
implementation or a new \texttt{scipy.sparse.linalg.eigsh} dependency.

Deliberately positioned as a \emph{cruder, opt-in} alternative, not a
replacement: it has no momentum label at all (a single scalar, not a
dispersion), does not require $D=1$, and -- critically -- never lets
\texttt{HL}/\texttt{HR} relax for whatever the second local eigenstate
actually represents, since they are exactly the ground state's own
converged environments. Measured directly against the two cases the
tangent-space ansatz was itself validated against: for the
field-polarized XX chain ($D=1$, exact answer known), it comes out
$\sim$10\% too high (5.5 vs.\ the exact 5.0) -- confirming it is a
genuinely cruder approximation, not a bug; for a genuinely entangled
($D>1$) dimerized Heisenberg chain -- exactly the case
\texttt{excitation\_gap} cannot handle -- it lands within $\sim$0.5\% of a
large finite open chain's own ED gap (\texttt{Spin\_Chain.get\_gap(mode=
"ED")} at \texttt{n\_sites=12,14,16}, the same finite-size-extrapolation
cross-check style \texttt{test\_n\_uc2\_dimerized\_chain\_matches\_finite\_
size\_extrapolation} already uses for the ground-state energy density), a
reassuring result precisely where the tangent-space ansatz offers nothing
at all. See \texttt{examples/idmrg/local\_excitation\_gap/main.py} for the
worked comparison.

\textbf{Tightening the local superblock gap further: \texttt{window=}
(\texttt{pyitensor/idmrg.py::local\_excitation\_gap\_windowed})}: rather
than adding Krylov vectors to the same frozen 2-site block --
\texttt{local\_excitation\_gap}'s own $\dim>3$ branch already
diagonalizes that exactly via deflated Lanczos, so more iterations there
cannot change the answer -- this grows the local block itself by
\texttt{window} extra \emph{free} physical sites on each side, using fresh
copies of the periodic per-sublattice MPO tensor
(\texttt{\_build\_automaton}) threaded onto \texttt{HL}/\texttt{HR}'s own
mpo axis via the same \texttt{\_relabel\_pos}/\texttt{\_fresh\_physical\_
copy} machinery \texttt{idmrg\_ground\_state}'s own growing loop uses
(recovering \texttt{HL}/\texttt{HR}'s mpo-axis Index by elimination from
their own 3 legs, since \texttt{local\_superblock} does not store it
directly). Both the ground state and the deflated first excited state are
then re-solved fresh within this larger block via
\texttt{\_lanczos\_ground\_state}, rather than reusing the growing
algorithm's own ground vector -- \texttt{window=0} reduces to exactly the
same effective Hamiltonian \texttt{local\_excitation\_gap} diagonalizes
(used as an internal consistency check on the construction, confirmed to
agree to Lanczos precision). Costs \texttt{d**(2*window)} more local
Hilbert space dimension per extra site pair (\texttt{d} = the physical
dimension), and is supported only for \texttt{n\_uc=1} -- widening needs
to know which sublattice position each extra inserted site takes, which
is not tracked for \texttt{n\_uc=2}.

Measured directly, this is a real, converging refinement, not just noise:
on the field-polarized XX chain ($D=1$, exact answer known) the error
drops monotonically from 10\% at \texttt{window=0} to
3.8\%/2.0\%/0.81\%/0.44\% at \texttt{window=1/2/4/6}; on a gapped,
genuinely entangled ($D>1$) transverse-field Ising chain,
\texttt{window=3} (an 8-site local block, gap 2.047) matches an 18-site
open finite chain's own ED gap (2.049) to $<$1\%, converging at least as
fast as growing the finite chain itself does ($n$=10/12/14/16/18 give ED
gaps 2.133/2.099/2.076/2.060/2.049). Reproducible to $\sim$1e-9 across
repeated calls despite the randomized Lanczos start, exactly like
\texttt{local\_excitation\_gap} itself. The convergence \emph{rate},
however, depends on the model's correlation length, not just its gap: on
the \texttt{S=1} Heisenberg chain (the Haldane gap, correlation length
$\sim$6 sites), \texttt{window=0/1/2} only move the estimate from
$\sim$29\%/24\%/21\% too high -- still improving, but far more slowly,
since a handful of extra sites barely dents a correlation length that
long, and the larger physical dimension ($d=3$ vs.\ $d=2$) makes each
extra site pair 9$\times$ more expensive rather than 4$\times$, making
large \texttt{window} impractical there. See
\texttt{examples/idmrg/local\_excitation\_gap/main.py} for the worked
comparison across models.

\textbf{Dynamical correlators, via a finite-window reduction to the
existing finite-chain KPM stack (\texttt{infinitechain.py}, not
\texttt{pyitensor/idmrg.py})}: unlike the static-correlator machinery
above, \texttt{Infinite\_Many\_Body\_Chain.kpm\_finite} deliberately
does \textbf{not} add any new Chebyshev-recursion code to
\texttt{idmrg.py}. Named \texttt{kpm\_finite} (not
\texttt{get\_dynamical\_correlator}, the finite-chain method it wraps)
so the name itself flags that this is a finite-window approximation,
not the exact infinite-size method described further below. $H$ is
extensive/unbounded in the thermodynamic limit, so there is no
infinite-chain analogue of the transfer-matrix formalism
\texttt{vev}/\texttt{correlator} use -- a literal KPM expansion of the
full infinite $H$ has no meaning (the same reason \texttt{apply\_mpo}
restricts \texttt{W\_bulk} to \emph{bounded} periodic operators, see
above). Instead, \texttt{kpm\_finite} builds an ordinary finite,
open-boundary \texttt{Many\_Body\_Chain}
(\texttt{itensor\_version="python"}, \texttt{n\_window} repeats of the
unit cell) with a Hamiltonian tiled from \texttt{self.\_h\_intra}/
\texttt{self.\_h\_inter} (\texttt{\_window\_hamiltonian}/
\texttt{\_shift\_terms}, new module-level helpers in
\texttt{infinitechain.py} -- pure term-list arithmetic, reusing the
same \texttt{MultiOperator.op} 0-based-site-index format
\texttt{\_canonicalize\_hamiltonian} already produces, just
concatenated across cells with an increasing site-index shift), places
the two operators at the window's central cell, and calls
\texttt{kpmdmrg.get\_dynamical\_correlator} on that temporary chain
directly -- the exact same code path (and, transitively,
\texttt{pyitensor/chain.py}'s \texttt{Chain.kpm\_dynamical\_correlator}/
\texttt{\_kpm\_moments\_full}/\texttt{\_scaled\_hamiltonian}) an
ordinary finite \texttt{Spin\_Chain}/\texttt{Fermionic\_Chain} already
uses, completely unmodified. \texttt{kpmdmrg.get\_dynamical\_correlator}
is called directly rather than through the usual public
\texttt{Many\_Body\_Chain.get\_dynamical\_correlator}/\texttt{dynamics.
get\_dynamical\_correlator} dispatch, to sidestep \texttt{dynamics.py}'s
own \texttt{is\_hermitian()} gate -- confirmed (this is
\texttt{set\_hamiltonian}'s own already-documented gap, not a new
finding) that \texttt{is\_hermitian()}'s \texttt{simplify()} step
false-rejects an ordinary cross-site
\texttt{Sx[i]*Sx[j]+Sy[i]*Sy[j]+Sz[i]*Sz[j]}-style term as
non-Hermitian, which would otherwise misroute a perfectly ordinary
Heisenberg-type window Hamiltonian to the non-Hermitian KPM/dynamics
path.

Correctness of the term-tiling itself (as opposed to the pre-existing,
independently-tested KPM math it feeds into) is checked in
\texttt{tests/test\_infinite\_chain.py} by an \emph{exact} ED
ground-state-energy comparison between \texttt{\_window\_hamiltonian}'s
output and an independently, by-hand-built equivalent finite chain --
the same finite-size-extrapolation pattern
\texttt{test\_n\_uc2\_dimerized\_chain\_matches\_finite\_size\_extrapolation}
already uses for \texttt{gs\_energy}, just applied to the Hamiltonian
construction directly instead of to iDMRG's own converged energy
density.

This is a genuine finite-size \textbf{approximation}, not an exact
infinite-size dynamical-correlator method -- see
\texttt{kpm\_finite}'s own docstring (and
\texttt{docs/user\_guide.tex}'s corresponding section) for the
light-cone/moment-count argument for why a fixed \texttt{delta} needs
\texttt{n\_window} chosen large enough (and, past a point,
\texttt{delta} itself coarsened) to avoid open-boundary reflections
contaminating the result, and
\texttt{examples/idmrg/dynamical\_correlator\_finite\_window/main.py}
for a worked \texttt{n\_window}-convergence sweep.

\textbf{A genuinely infinite-chain KPM is possible in principle, but
not implemented here.} Two separate ingredients would be needed, both
reusing machinery this module already has, not a small patch on top of
\texttt{kpm\_finite}:

\begin{enumerate}
\item Cap the window's two ends with the \emph{converged environment}
-- the \texttt{HL}/\texttt{HR} accumulated Hamiltonian blocks
\texttt{idmrg\_ground\_state} already builds during growth, or an
equivalent construction from the fixed-point machinery
\texttt{apply\_mpo}/\texttt{imps\_overlap} already use -- instead of
plain open boundaries: infinite boundary conditions (IBC,
Phien/McCulloch/Vidal, PRB 86, 245107 (2012)). This removes a real
error source \texttt{n\_window} alone cannot fix no matter how large:
an open chain's own ground state carries boundary artifacts (e.g.
Friedel-oscillation-like features) that contaminate even the
\emph{central} region, not just the two edges, because the window's
own DMRG ground-state search has no way to know it should match the
true infinite bulk rather than terminate at a physical edge.
\item Independently, let the \emph{active} window grow by roughly one
site per Chebyshev moment as the recursion proceeds (reusing the
converged unit-cell tensors to extend the state on demand, similar in
spirit to a ``growing window'' real-time TEBD/TDVP simulation on an
infinite chain), rather than fixing a static \texttt{n\_window} up
front. This removes the \texttt{n\_window}-vs-moment-count constraint
entirely, bounding cost only by the usual bond-dimension/computational
budget of any DMRG calculation, not by a fixed geometric window
guessed in advance.
\end{enumerate}

Assembling these into a working IBC-window KPM implementation is a
new, nontrivial feature in its own right -- flagged as a known,
deliberate follow-up rather than attempted here (see
\texttt{kpm\_finite}'s own docstring for the same note, addressed to a
caller rather than a maintainer).

\textbf{Infinite-boundary-condition (IBC) window construction and
real-time dynamical correlator
(\texttt{pyitensor/idmrg\_window.py})}: following Milsted/Vanderstraeten
et al., ``Infinite boundary conditions for response functions and limit
cycles in iDMRG'' (arXiv:1804.09163), Sec.~V.1 -- the item-1 ingredient
flagged just above as unimplemented future work. The paper's method
time-evolves a finite window of tensors around a local perturbation,
with everything outside the window frozen at the uniform ground state
and capped by ``the Hamiltonian terms outside the window, projected
onto the reduced $D$-dimensional Hilbert space of the left/right
block'' -- which turns out to be \emph{exactly} what
\texttt{idmrg\_ground\_state}'s own growth loop already computes as
\texttt{HL}/\texttt{HR}, just discarded once converged.
\texttt{idmrg.py} was extended (\texttt{IDMRGResult}'s own
\texttt{env\_HL}/\texttt{env\_HL\_bra}/\texttt{env\_HL\_ket}/
\texttt{env\_HL\_mpo} and mirrored \texttt{env\_HR\_*} fields, plus
\texttt{W\_bulk}) to snapshot these -- specifically, \texttt{HL}/
\texttt{HR} as they stood \emph{entering} the last executed
macro-iteration (captured via a per-macro-iteration
\texttt{env\_window\_boundary} snapshot, overwritten every iteration so
whatever remains at loop exit is the one consistent with the
\emph{returned} \texttt{U\_list}'s own edge bonds: \texttt{U\_list[0]}'s
left bond is exactly that snapshot's \texttt{HL\_ket},
\texttt{U\_list[n\_uc-1]}'s right bond exactly its \texttt{HR\_ket}, by
construction of how \texttt{svd()} splits the local 2-site solve) -- no
new environment solve needed, only exposing what already exists.

\texttt{idmrg\_window.py}'s \texttt{build\_window(result, n\_window)}
tiles \texttt{n\_window} periodic repeats of \texttt{U\_list} (the
window's MPS) and \texttt{W\_bulk} (its MPO) into an ordinary
\texttt{mpscontainer.MPS}/\texttt{MPO} pair (\texttt{\_tile\_periodic},
fresh Link Indices at every copy boundary to avoid the
identity-collision hazard \texttt{\_relabel\_pos}/\texttt{\_project\_
channel} already warn about elsewhere in this module), with the very
first/last tensors' outer legs left as \texttt{env\_HL\_ket}/
\texttt{env\_HL\_mpo}/\texttt{env\_HR\_ket}/\texttt{env\_HR\_mpo} so
they attach directly to the environment caps. Environment extension
\emph{through} the window deliberately reuses \texttt{idmrg.py}'s own
explicit-Index \texttt{\_extend\_HL}/\texttt{\_extend\_HR}
(\texttt{\_extend\_through\_left}/\texttt{\_extend\_through\_right}),
not \texttt{dmrg.py}'s generic \texttt{\_Chain}/\texttt{\_link\_at}-
based \texttt{\_extend\_left}/\texttt{\_extend\_right}: the latter
infers a site's left/right Link purely by looking for a same-Index
neighbor \emph{within the chain itself}, correct for an ordinary
finite chain (whose boundary sites truly have no further leg, see
\texttt{mpscontainer.py}'s own docstring) but wrong here, since the
window's own edge sites \emph{do} carry one more leg each (connecting
to \texttt{env\_HL}/\texttt{env\_HR}, tensors living outside the
\texttt{\_Chain} object entirely) that \texttt{\_link\_at} cannot see --
confirmed directly: using the generic functions silently
mis-contracted the bra/ket legs at the window's own edge sites, since
\texttt{\_link\_at}'s neighbor lookup returns ``no such leg'' for a
boundary site regardless of whether the tensor itself actually has one.

\textbf{A confirmed, non-obvious normalization finding}:
\texttt{window\_total\_energy} (the capped window's total energy, via
\texttt{\_extend\_through\_left} then closing against
\texttt{env\_HR}) is \emph{not} simply \texttt{e0 * (site count)} --
plugging in the converged \texttt{U\_list} tensors with both
\texttt{env\_HL\_ket} and \texttt{env\_HR\_ket} left as free/dangling
legs (rather than contracted against a specific boundary vector) gives
$\langle\text{window}|\text{window}\rangle = \dim(\texttt{env\_HR\_ket})$,
not $1$ (an isometry-chain argument: each site's own left-canonical
$U$ collapses the sum over every index except the very last dangling
one, leaving that dimension as the Frobenius norm-squared -- confirmed
numerically against an independent Lanczos-Rayleigh-quotient
computation using the identical \texttt{env\_HL}/\texttt{env\_HR}/
\texttt{W\_pL}/\texttt{W\_pR}, to machine precision) --
\texttt{window\_total\_energy} divides by it. Even normalized, the
\emph{absolute} total still carries a large, \texttt{n\_window}-
independent baseline from whatever \texttt{env\_HL}/\texttt{env\_HR}
themselves represent (every macro-iteration already absorbed before
the snapshot) -- so \texttt{window\_energy\_density(result,
n\_window)}, a finite difference between window sizes
\texttt{n\_window} and \texttt{n\_window+1} divided by \texttt{n\_uc}
(mirroring \texttt{idmrg\_ground\_state}'s own
\texttt{density = (energy-prev\_energy)/(2*n\_uc)} diagnostic), is the
well-posed sanity check this module validates against, not
\texttt{window\_total\_energy} directly. \textbf{A second, real bug
this same finite-difference check caught}: an earlier version compared
window sizes \texttt{n\_window} and \texttt{n\_window+n\_uc} (not
\texttt{n\_window+1}) while still dividing by \texttt{n\_uc} --
\texttt{n\_window} counts \emph{unit-cell copies}, so \texttt{+n\_uc}
copies adds \texttt{n\_uc*n\_uc} physical sites, not \texttt{n\_uc} --
invisible for \texttt{n\_uc=1} (where \texttt{n\_uc*n\_uc==n\_uc}), but
confirmed directly to inflate every \texttt{n\_uc=2} result by
\emph{exactly} a factor of 2 for an easy, cleanly-converged test
model (an exact ratio, not convergence noise, which is what exposed it
as a real off-by-\texttt{n\_uc} bug rather than an iDMRG convergence
limitation).

\textbf{Real-time evolution, shifted overlaps, and $S(k,\omega)$}:
\texttt{window\_tdvp\_step} two-site-TDVP-evolves a capped window in
place (mirroring \texttt{tdvp.py}'s own \texttt{tdvp\_step}), but
\texttt{tdvp.py}'s own sweep functions cannot be reused directly --
they build each bond's local effective Hamiltonian and environments via
\texttt{mpsalgebra.py}'s \texttt{\_link\_at}, which finds a site's
left/right Link purely by looking for a same-Index \emph{neighbor in
the chain itself}; correct for an ordinary finite chain (whose boundary
sites truly have no further leg) but wrong for a window, whose edge
sites \emph{do} carry one more leg each (connecting to
\texttt{env\_HL}/\texttt{env\_HR}, tensors outside the chain object
entirely) that \texttt{\_link\_at} cannot see -- confirmed directly,
passing a window straight through \texttt{tdvp.py}'s own sweep
functions silently mis-contracts the bra/ket legs at the window's own
edge sites. \texttt{idmrg\_window.py}'s own
\texttt{\_window\_two\_site\_heff}/\texttt{\_window\_one\_site\_heff}
read a site's left/right Link directly off its own tensor instead
(always present for a window, unlike an ordinary chain's boundary), and
\texttt{\_all\_left\_environments\_window}/
\texttt{\_all\_right\_environments\_window} build environments via
idmrg.py's own explicit-Index \texttt{\_extend\_HL}/\texttt{\_extend\_HR}.
\texttt{apply\_local\_operator}/\texttt{local\_expectation} add the
perturbation and readout primitives; \texttt{local\_expectation} needed
its own real fix, distinct from the TDVP one: a plain observable's own
boundary weighting is \emph{not} symmetric between the window's two
ends, because \texttt{U\_list} is left-canonical (its own
left-canonicality condition \emph{is} the statement that the left
transfer-matrix fixed point is trivially the identity, needing no
weighting) -- the \emph{right} boundary needs the dominant right fixed
point (\texttt{\_all\_right\_fixed\_points}), exactly like
\texttt{onsite\_expectation}/\texttt{two\_point\_correlator} already
rely on; an earlier version used a bare trace on both ends and produced
a visibly non-uniform profile for a converged, translation-invariant
ground state that should read exactly uniform everywhere.

\texttt{dynamical\_correlator\_td}/\texttt{snapshot\_correlator}
reconstruct $S(x,t)=\langle\psi|A_x e^{-iHt}B_0|\psi\rangle$ for
\emph{every} $x$ from a \emph{single} window evolution (the paper's own
headline efficiency result), by inserting operator $A$ directly into
the \emph{bra} side of the overlap at the shifted absolute position,
padding whichever side's explicit tensor list falls short with
unevolved \texttt{U\_list} copies (\texttt{\_padded\_arrays}) -- valid
because, outside each window's own causal cone, its tensors already
equal \texttt{U\_list} exactly. This is a simplification of the paper's
own Eq.~7 to $t_1=0$ (the naive Eq.~3 form) rather than the full
two-branch trick (evolving a \emph{second} window backward by $t_1$
too, doubling the accessible total time for the same TDVP cost) --
\texttt{window\_tdvp\_step} already supports backward evolution via a
negative \texttt{dt}, so only the second branch and its overlap
bookkeeping are missing, a documented follow-up. \texttt{connected=True}
(the default) subtracts the disconnected background $\langle A\rangle
\langle B\rangle$ (\texttt{onsite\_expectation}) -- confirmed directly
to matter: the \emph{raw} correlator approaches $\langle A\rangle
\langle B\rangle$ (not $0$) at large $|x|$, and summing that
non-decaying background over $x$ with $e^{-ikx}$ produces a spurious,
dominant $k=0$ contribution with no discernible dispersion.
\texttt{dynamical\_correlator\_komega} does the spatial DFT
($S(k,t)=\sum_x e^{-ikx}S(x,t)$) then reuses \texttt{timedependent.py}'s
own \texttt{\_fourier\_transform\_correlator} unchanged (it was
explicitly factored out for other time-domain submodes to reuse, see
\texttt{tdz.py}'s own use of it), so \texttt{delta} means the same
Lorentzian-broadening thing here as in every other dynamical-correlator
submode in this codebase.

Validated: $S(x,t{=}0)$ matches idmrg.py's own exact static
\texttt{two\_point\_correlator} to machine precision ($10^{-16}$) for
every $x$ tested (\texttt{tests/test\_idmrg\_window.py}); $S(k,\omega)$
(with \texttt{connected=True}) shows a clear, $k$-dependent peak
instead of a featureless background; cross-checked against
\texttt{kpm\_finite} (an independent approximation scheme --
open-boundary window + KPM, vs.\ this method's IBC window + TDVP) for
the same physical correlator -- an exact match isn't expected, but both
agree on where the dominant spectral weight sits
(\texttt{tests/test\_infinite\_chain.py},
\texttt{examples/idmrg/td\_dynamical\_correlator/main.py}).

\texttt{Infinite\_Many\_Body\_Chain.td\_dynamical\_correlator}
(\texttt{infinitechain.py}) wires this up as the public API, mirroring
\texttt{kpm\_finite}'s own docstring conventions (\texttt{p\_i}
sublattice position, an \texttt{n\_window} convergence caveat) --
\texttt{itensor\_version="python"} only (raises
\texttt{NotImplementedError} otherwise), calls \texttt{gs\_energy()}
automatically like \texttt{vev}/\texttt{correlator}. See
\texttt{docs/user\_guide.tex}'s own iDMRG section for the user-facing
version.

\subsection{Operator representation: \texttt{MultiOperator}}

Hamiltonians and observables are built as \texttt{MultiOperator} objects
(\texttt{multioperator.py}) --- sums of products of named single-site
operators, e.g.\ \texttt{Sx[i]*Sx[j]}. This is a backend-agnostic
intermediate representation: the \emph{same} \texttt{MultiOperator} is
later either

\begin{itemize}
  \item converted into a plain list of
    \texttt{(coefficient, [(opname, site), \ldots])} terms and written
    directly into an ITensor \texttt{AutoMPO}/\texttt{HTerm} by the
    C++/pyitensor DMRG backends, or
  \item converted into a sparse matrix via
    \texttt{multioperator.MO2matrix} for the ED backend
    (\texttt{edtk/edchain.py}).
\end{itemize}

\texttt{multioperatortk/} holds the supporting machinery: Jordan-Wigner
string threading for fermionic operators, static/long-range operator
construction, and sympy-based symbolic building. Spinless fermionic
operators (\texttt{C}/\texttt{Cdag}, one fermionic mode per site) are
threaded by \texttt{jordanwigner.py}; flavor-resolved operators on a
native spinful site
(\texttt{Cup}/\texttt{Cdn}/\texttt{Cdagup}/\texttt{Cdagdn}, two
fermionic modes per site, \texttt{Spinful\_Fermionic\_Chain\_Native})
are threaded by the separate \texttt{jordanwigner\_spinful.py} instead,
dispatched by operator name in \texttt{multioperator.jordan\_wigner()}.
The two never mix within one term (the operator name sets are
disjoint), and \texttt{jordanwigner\_spinful.py} always uses the generic
per-factor dressing recipe (no separate optimized 2-/4-point path the
way \texttt{jordanwigner.py} has for plain \texttt{C}/\texttt{Cdag}),
since that recipe is already exact for an arbitrary product/order of
factors --- see its module docstring.

\subsection{Backend dispatch}

\texttt{mode.py} decides, per call, whether a calculation runs via DMRG
or ED (the \texttt{get\_mode}/\texttt{run} functions): DMRG unless
\texttt{self.mode} forces ED, or unless \texttt{self.itensor\_version} is
\texttt{2} or \texttt{3} and the corresponding compiled pybind11
extension isn't available (\texttt{cppext.available(version)}) --- in
which case DMRGPY silently falls back to ED.
\texttt{itensor\_version="python"} never falls back this way, since it
has no compiled-extension precondition at all. Most public
\texttt{Many\_Body\_Chain} methods accept \texttt{mode="DMRG"|"ED"} so
results can be cross-validated between the two solver families.

\begin{center}
\begin{tabular}{@{}lp{0.32\textwidth}p{0.18\textwidth}l@{}}
\toprule
\texttt{itensor\_version} & Engine & Requires & Fallback \\
\midrule
\texttt{2} & ITensor v2, in-process C++ (mpscpp2) & compiled pybind11 ext. & ED \\
\texttt{3} (default) & ITensor v3, in-process C++ (mpscpp3) & compiled pybind11 ext. & ED \\
\texttt{"python"} & pure-Python \texttt{pyitensor/} & NumPy/SciPy only & none \\
\texttt{"julia\_live"} & live in-process Julia (mpsjulialive), ITensors.jl & pyjulia + Julia & none (per-feature) \\
\bottomrule
\end{tabular}
\end{center}

Regardless of \texttt{itensor\_version}, if \texttt{self.mode} is forced
to \texttt{"ED"}, or the requested backend isn't available, calculations
run through \texttt{edtk/edchain.py} (\texttt{EDchain}): dense/sparse
operators built directly in Python/NumPy/SciPy (\texttt{pyfermion/},
\texttt{pyspin/}, \texttt{pyboson/}, \texttt{pyzn/} provide per-statistics
many-body operator construction; \texttt{edtk/one2many.py} promotes
single-site operators to the full Hilbert space), diagonalized with
\texttt{scipy.sparse.linalg}.

\subsection{The in-process C++ backend (ITensor v2 and v3)}

\texttt{mpscpp2/bindings.cc} and \texttt{mpscpp3/bindings.cc} each
compile a pybind11 extension (\texttt{mpscpp2/\_dmrgcpp*.so} /
\texttt{mpscpp3/\_dmrgcpp*.so}) that runs DMRG entirely in-process against
its own vendored ITensor copy --- no task files, no operator files, no
subprocess, no temp directory. \texttt{mpscpp3} is a line-by-line port of
\texttt{mpscpp2} to the ITensor v3 API (which merges
\texttt{IQTensor}/\texttt{IQIndex} into a single \texttt{ITensor}/
\texttt{Index} type carrying QN blocks on the Index), exposing an
identical \texttt{Chain} class and Python surface. \texttt{cppext.py}
lazily imports whichever compiled extension a chain needs
(\texttt{get\_backend(version)}, \texttt{available(version)}).

\texttt{mpscppN/chain\_session.h}'s \texttt{Chain} class is the stateful,
in-process session --- one instance per \texttt{Many\_Body\_Chain} (held
as \texttt{self.\_session}) --- with a method per DMRG task:
\texttt{gs\_energy}, \texttt{vev}, \texttt{apply\_operator}, KPM/CVM
dynamical correlators, \texttt{correlation\_matrix}, \texttt{reduced\_dm},
\texttt{excited\_states}, time evolution, and more.

The session caches its ground-state energy and (for KPM) both band
edges, but \texttt{Chain::set\_hamiltonian} unconditionally invalidates
those caches --- so the Python side
(\texttt{groundstate.py::gs\_energy\_single}) only re-sends the
Hamiltonian when its \texttt{to\_terms()} output, any solver parameter
a re-run would pick up (\texttt{maxm}, \texttt{nsweeps},
\texttt{cutoff}, \texttt{noise}, the effective MPO bond dimension), or
the session object itself actually changed since the last send
(\texttt{\_session\_ham\_cache}). Repeated calculations
on an unchanged Hamiltonian (e.g.\ successive
\texttt{get\_dynamical\_correlator} calls, each of which re-verifies
the ground state) then hit the session's caches instead of re-running
warm DMRG sweeps and band-edge solves. Code paths that \emph{want} a
fresh solve of the same Hamiltonian either pass through
\texttt{restart()}/\texttt{set\_hamiltonian} (which force DMRG via
\texttt{skip\_dmrg\_gs=False}) or clear \texttt{\_session\_ham\_cache}
explicitly (see \texttt{groundstate.py}'s best-of-$n$ loops).

Notable, deliberate implementation details (not bugs to ``fix''):

\begin{itemize}
  \item \texttt{mpscpp3} builds every site with \texttt{ConserveQNs=false}
    and starts DMRG from an actual \texttt{randomMPS}, not a plain
    product state --- matching v2's real (unconstrained-search) behavior
    rather than ITensor v3's stricter, QN-conserving-from-the-start
    convention, at the cost of losing the QN block-sparsity speedup for
    \texttt{itensor\_version=3}.
  \item \texttt{mpscpp3}'s (and \texttt{pyitensor}'s) real-time MPS
    evolution defaults to a proper two-site TDVP integrator (vendored
    under \texttt{mpscpp3/TDVP/}; \texttt{pyitensor/tdvp.py} for the
    pure-Python backend), selectable via
    \texttt{Many\_Body\_Chain.tevol\_method} (default \texttt{"TDVP"},
    \texttt{itensor\_version} 3 or \texttt{"python"} only);
    \texttt{mpscpp2} has no TDVP and always uses a hand-rolled 2nd-order
    Taylor expansion of $\exp(-i\,dt\,H)$ as an MPO instead.
  \item \texttt{tevol\_method="TDVP\_GSE"} (\texttt{itensor\_version} 3 or
    \texttt{"python"} only, same support as plain \texttt{"TDVP"}) runs
    one-site TDVP with Krylov \emph{global subspace expansion} (GSE)
    beforehand for the first \texttt{tdvp\_gse\_sweeps} steps (default 3)
    --- the scheme of Yang \& White, arXiv:2005.06104/Phys.\ Rev.\ B 102,
    094315 (2020): a Krylov subspace $\{\psi,H\psi,H^2\psi,\dots\}$ of
    dimension \texttt{tdvp\_gse\_krylov\_order} (built via repeated MPO
    application) enlarges the MPS's local bond bases \emph{without
    changing the represented state}, giving one-site TDVP (which
    conserves bond dimension exactly on its own, unlike two-site) room to
    grow into the entanglement the subsequent evolution generates.
    \texttt{itensor\_version=3}
    (\texttt{Chain::global\_subspace\_expand()}/
    \texttt{Chain::tdvp\_step(...,num\_center=1)},
    \texttt{mpscpp3/TDVP/basisextension.h}, vendored unmodified from
    upstream) and \texttt{"python"} (\texttt{pyitensor/gse.py}) implement
    this. A v2-API port (\texttt{mpscpp2/TDVP/}) was attempted and
    briefly landed: numerically correct (verified against ED and against
    v3/\texttt{"python"}, exact agreement on a 6-site cross-check), built
    on a from-scratch Lanczos \texttt{applyExp} (v2's own
    \texttt{itensor/iterativesolvers.h} never had one) since v2's stock
    \texttt{LocalMPO} also turned out to have no working one-site local
    Hamiltonian at all (confirmed in \texttt{LocalMPO::position()}, which
    unconditionally builds a two-site block regardless of
    \texttt{numCenter()}), paired with two-site TDVP instead as a result.
    It was reverted after a severe, unresolved performance regression at
    $n\gtrsim10$ sites (the dynamical-correlator step didn't finish in 25
    minutes at $n=12$, versus under a second for the same computation on
    \texttt{itensor\_version=3}) that couldn't be root-caused in the time
    available --- see git history around \texttt{mpscpp2/TDVP/} if
    picking this up again.
  \item All three session backends implement a real non-Hermitian DMRG
    (\texttt{Chain::nhdmrg}, driven by \texttt{nhdmrg.py}, exposed as
    \texttt{Many\_Body\_Chain.nhdmrg()}): a port of ITensorNHDMRG.jl's
    ``onesided''+``fidelity'' algorithm that optimizes a biorthogonal
    left/right eigenpair of a non-Hermitian $H$, targeting the eigenvalue
    with smallest real part. \texttt{mpscpp3/chain\_session.h}'s copy is
    the annotated original (including two deliberate deviations from the
    reference for Re-degenerate spectra); \texttt{mpscpp2}'s is its
    v2-API back-port, and \texttt{pyitensor/nhdmrg.py} is the pure-Python
    port (built on \texttt{dmrg.py}'s environment/matvec machinery;
    unlike \texttt{dmrg.py} it \emph{does} implement the noise term,
    because for NH-DMRG it measurably matters). \texttt{groundstate.py}'s
    non-Hermitian \texttt{gs\_energy} branch routes to it for
    \texttt{itensor\_version} 2, 3 and \texttt{"python"}; the MPS Arnoldi
    route (default \texttt{algebra/arpacktk.py}, an Implicitly Restarted
    Arnoldi Method ported from ARPACK; \texttt{algebra/arnolditk.py}'s
    explicit-restart Arnoldi remains available for comparison, see
    \texttt{mpsalgebra.py}) remains as the fallback for other backends.
    The adjoint MPO is built from
    \texttt{MultiOperator.get\_dagger()}'s terms on the Python side, and
    since the non-Hermitian energy is not variational, \texttt{nhdmrg.py}
    certifies each run by the eigen-residual and redraws a fresh random
    start when a run stalls.
  \item A few pre-existing bugs in the original (pre-refactor) file-based
    backend are deliberately reproduced rather than silently fixed --- see
    the call-site comments in \texttt{chain\_session.h} for both
    versions.
  \item \texttt{itensor\_version} 2, 3 (and \texttt{"python"}) expose MPO
    algebra directly on already-built operators: \texttt{StaticOperator}
    supports \texttt{+}, \texttt{-}, unary \texttt{-}, and scalar
    \texttt{*}/\texttt{/} between two \texttt{StaticOperator}s (or a
    scalar), on top of the pre-existing \texttt{*} for
    \texttt{StaticOperator*MPS}/\texttt{StaticOperator} contraction.
    \texttt{A + B} (\texttt{Chain::sum\_operators}, a public wrapper each
    backend already had internally as a private \texttt{sum\_mpo()}
    helper used by \texttt{custom\_exp()}/\texttt{evoloperator()}) is a
    compressed direct sum at the tensor-network level --- algorithmically
    the same construction as ITensorMPS.jl's \texttt{+(::MPO, ::MPO)}
    (\texttt{abstractmps.jl}'s default \texttt{"densitymatrix"} algorithm)
    --- not a MultiOperator-level symbolic sum (which already existed, see
    \S4.2, and remains the preferred way to combine Hamiltonians before
    ever building an MPO). It exists for combining operators that only
    exist as already-built \texttt{StaticOperator}s, e.g.\ two
    independently constructed products or exponentials.
    \texttt{"julia\_live"} doesn't implement this yet
    (\texttt{StaticOperator.\_\_add\_\_} raises
    \texttt{NotImplementedError} rather than silently doing something
    backend-specific).
\item \texttt{mpscpp3}-specific: \texttt{Chain::gs\_energy\_generalized}
    (\texttt{mpscpp3/chain\_session.h}, exposed as
    \texttt{Many\_Body\_Chain.gs\_energy\_generalized}) solves the
    generalized eigenproblem $H|\psi\rangle=\lambda A|\psi\rangle$ for a
    Hermitian positive-definite metric MPO $A$ via a self-consistent
    Lagrange-multiplier iteration: each outer iteration builds the MPO
    $H-\lambda A$ (\lstinline|sum()|) from the current $\lambda$
    estimate, sweeps it with ITensor v3's own \lstinline|dmrg()| for a
    single length-1 \lstinline|Sweeps| step (looping outer iterations by
    hand rather than handing the whole schedule to one \lstinline|dmrg()|
    call, since $\lambda$ must be updated between sweeps --- noise is
    halved off partway through exactly like \texttt{nhdmrg()}'s own
    per-sweep loop above), then updates $\lambda$ to the swept state's
    generalized Rayleigh quotient
    $\langle\psi|H|\psi\rangle/\langle\psi|A|\psi\rangle$. A
    line-for-line port of \texttt{pyitensor/dmrg.py}'s
    \texttt{dmrg\_generalized} (\S4.5) against real ITensor v3 calls
    instead of the hand-rolled two-site sweep pyitensor uses;
    \texttt{mpscpp2} has no analogous session method yet, so
    \texttt{itensor\_version=2} still raises
    \texttt{NotImplementedError}. See
    \texttt{examples/groundstate/dmrg\_generalized\_benchmark}, which
    heads all three available routes (pyitensor, v3, and
    \texttt{algebra/arpacktk.py}'s pre-existing ARPACK-mode-2
    \texttt{mpsiram\_generalized}) up against each other on the same
    test problem --- v3 is consistently the fastest of the three at the
    (small) sizes benchmarked there.
\item \texttt{mpscpp3}-specific: \texttt{Chain::nhdmrg\_generalized}
    (exposed as \texttt{Many\_Body\_Chain.gs\_energy\_generalized} for a
    non-Hermitian \texttt{self.hamiltonian}) is the non-Hermitian
    counterpart of the bullet above, generalizing \texttt{Chain::nhdmrg}
    (\S4.4's own NH-DMRG bullet) the same way
    \texttt{gs\_energy\_generalized} generalizes plain
    \texttt{gs\_energy()}. Since the metric $A$ is Hermitian,
    $(\lambda A)^\dagger=\bar\lambda A$, so the same self-consistent
    Lagrange-multiplier trick carries over with a \emph{complex} $\lambda$
    and \emph{biorthogonal} (not plain) expectation values: each outer
    iteration shifts both $H$ and its precomputed adjoint $H^\dagger$ by
    $\lambda A$/$\bar\lambda A$ respectively, runs one ordinary NH-DMRG
    sweep (a new private \texttt{nhdmrg\_one\_sweep} helper, factored out
    of \texttt{Chain::nhdmrg}'s own loop body) against the shifted pair,
    then updates $\lambda$ to
    $\langle\psi_L|H|\psi_R\rangle/\langle\psi_L|A|\psi_R\rangle$.
    Because \texttt{nhdmrg\_one\_sweep}'s two-site solve is hand-rolled
    directly against \texttt{arnoldi\_smallest\_real}/manual ITensor
    contractions rather than a call into ITensor v3's own
    \texttt{dmrg()}, it does \emph{not} need
    \texttt{gs\_energy\_generalized}'s own short-chain guard (confirmed
    directly, a 2-site chain runs it without aborting) --- a real
    asymmetry between the two \texttt{mpscpp3}-specific bullets in this
    section, not an oversight. Pyitensor gained the identical algorithm
    first (\texttt{nhdmrg.py::nhdmrg\_generalized}); \texttt{mpscpp2}
    still has no analogous session method for either backend's
    non-Hermitian generalized solver.
\end{itemize}

\subsection{The pure-Python backend (\texttt{pyitensor/})}

\texttt{pyitensor/} is a from-scratch, pure-Python/NumPy/SciPy
reimplementation of exactly the ITensor v3 API subset
\texttt{mpscpp3/chain\_session.h} uses: an \texttt{Index}/\texttt{ITensor}
tensor core, the same ten site types, an \texttt{AutoMPO} term compiler
with its own Jordan-Wigner threading, MPS/MPO algebra, two-site
Lanczos-based DMRG (ground and excited states via the overlap-penalty
method), two-site Krylov-based TDVP, and a \texttt{chain.Chain} facade
exposing the identical method surface as the compiled backends'
\texttt{self.\_session}. Because of that shared surface, every call site
elsewhere in DMRGPY that accepts \texttt{itensor\_version in (2, 3)}
treats \texttt{"python"} as a third, always-available option with no
separate code path.

It exists so DMRG/TDVP work with zero compiler/pybind11 dependency, at
the cost of being slower than compiled ITensor by default (no
block-sparsity, no JIT) --- see Section~\ref{sec:perf} for how much
slower in practice, and how \texttt{numba}/\texttt{jax} narrow that gap.

One partial exception to that shared-surface rule:
\texttt{dmrg.py::dmrg\_generalized} (exposed as
\texttt{Chain.gs\_energy\_generalized}/
\texttt{Many\_Body\_Chain.gs\_energy\_generalized}) solves the generalized
eigenproblem $H|\psi\rangle=\lambda A|\psi\rangle$ for a Hermitian
positive-definite metric MPO $A$ --- a self-consistent Lagrange-multiplier
iteration built directly on top of \texttt{dmrg()}'s own per-sweep
machinery (factored into a shared \texttt{\_dmrg\_one\_sweep} helper):
each outer iteration rebuilds the MPO $H-\lambda A$ from the current
$\lambda$ estimate, sweeps it with the ordinary two-site solver, then
updates $\lambda$ to the swept state's generalized Rayleigh quotient
$\langle\psi|H|\psi\rangle/\langle\psi|A|\psi\rangle$. \texttt{mpscpp3}
has since gained a line-for-line port of this same algorithm against
real ITensor v3 calls (\S4.4's own \texttt{mpscpp3}-specific bullet), so
\texttt{itensor\_version} 3 supports it too now --- only
\texttt{mpscpp2} (\texttt{itensor\_version=2}) and
\texttt{"julia\_live"} still raise \texttt{NotImplementedError}, not
having an analogous session method yet.

\texttt{nhdmrg.py::nhdmrg\_generalized} is the non-Hermitian counterpart,
generalizing NH-DMRG (\S4.4's own bullet above, and pyitensor's own port
of it) exactly the way \texttt{dmrg\_generalized} generalizes plain
\texttt{dmrg()}: \texttt{groundstate.gs\_energy\_generalized} checks
\texttt{self.hamiltonian} for Hermiticity and transparently dispatches a
non-Hermitian one here instead of raising. Since $A$ is Hermitian,
$(\lambda A)^\dagger=\bar\lambda A$, so the same trick carries over with
a \emph{complex} $\lambda$ and \emph{biorthogonal} (not plain)
expectation values: each outer iteration shifts both $H$ and its
precomputed adjoint $H^\dagger$ by $\lambda A$/$\bar\lambda A$
respectively, runs one ordinary NH-DMRG sweep
(\texttt{\_nhdmrg\_one\_sweep}, factored out of \texttt{nhdmrg()}'s own
loop body the same way \texttt{\_dmrg\_one\_sweep} was) against the
shifted pair, then updates $\lambda$ to
$\langle\psi_L|H|\psi_R\rangle/\langle\psi_L|A|\psi_R\rangle$.
\texttt{mpscpp3} has since gained a line-for-line port of this algorithm
too (\S4.4's own second \texttt{mpscpp3}-specific bullet), so
\texttt{itensor\_version} 3 supports a non-Hermitian
\texttt{self.hamiltonian} here just like the Hermitian case above ---
\texttt{mpscpp2} is still the only backend without an analogous session
method for either. See
\texttt{examples/non\_hermitian/nhdmrg\_generalized\_benchmark}, which
cross-checks both implementations against \texttt{mpsiram\_generalized}
(ARPACK mode 2 needs no adaptation at all for a non-Hermitian primary
operator, only its own $M$ positive-definite precondition was ever
required) on a non-Hermitian test problem: the same accuracy/speed
pattern as the Hermitian benchmark holds against ARPACK, but unlike the
Hermitian case v3 isn't consistently faster than pyitensor here ---
NH-DMRG's per-bond cost already pays for \emph{two} Arnoldi solves (right
block and its adjoint) regardless of backend, narrowing the
compiled-vs-pure-Python gap relative to plain ground-state DMRG's single
local diagonalization per bond.

\subsection{The Julia backend (\texttt{mpsjulialive/})}

\texttt{itensor\_version="julia\_live"} drives a live, in-process Julia
session (via
\href{https://github.com/JuliaPy/PythonCall.jl}{\texttt{juliacall}/PythonCall.jl},
\texttt{mpsjulialive/juliasession.py}) with its own set of modules
mirroring the top-level ones (\texttt{mpsjulialive/groundstate.py},
\texttt{vev.py}, \texttt{mps.py}, \texttt{mpo.py}, \texttt{dynamics.py},
\ldots) that talk to Julia's ITensors.jl instead of the C++ pybind11
extension. This is an independent implementation, not a shared protocol
with the C++ path --- a feature missing on one side simply isn't
implemented there yet.

\texttt{juliacall} replaced an earlier PyJulia-based bridge: PyJulia
requires the Julia build's PyCall to be linked against a matching
libpython (the same class of ABI landmine documented in
Section~\ref{sec:install} for the C++ extension) and needs a
\texttt{compiled\_modules=False} workaround for slow/broken
precompilation. \texttt{juliacall} avoids both, and self-provisions Julia
plus the required packages (\texttt{ITensors}, \texttt{ITensorMPS},
\texttt{ITensorNHDMRG}, pinned in \texttt{src/dmrgpy/juliapkg.json}) into
its own managed project the first time \texttt{mpsjulialive} is imported
--- no manual Julia download step. One porting detail worth knowing:
unlike PyJulia, \texttt{juliacall} does \emph{not} implicitly convert a
Python \texttt{list} of \texttt{str} into a Julia \texttt{Vector\{String\}}
(it wraps it as a lazy \texttt{PyList\{Any\}}, which fails to dispatch
against strictly-typed Julia functions like \texttt{sites.jl}'s
\texttt{get\_sites}); every call site that used to rely on that now goes
through \texttt{juliasession.to\_julia\_strvec()}, which forces the
conversion explicitly via \texttt{juliacall.convert}.

The KPM dynamical correlator's Chebyshev-moment recursion runs natively
in Julia (\texttt{mpsjulialive/kpm.jl}'s \texttt{kpm\_moments\_full}/
\texttt{kpm\_moments\_accelerated}, driven through \texttt{kpm.jl}'s own
\texttt{apply\_op}/\texttt{mpsalgebra.jl}'s \texttt{summps}), one call
per correlator rather than one Python$\leftrightarrow$Julia round trip
per Chebyshev step. \texttt{mpsjulialive/dynamics.py} only builds the
scaled-Hamiltonian MPO and the two seed wavefunctions in Python
(mirroring \texttt{pyitensor/chain.py}'s own pure-Python KPM algorithm,
since \texttt{julia\_live} has no single \texttt{Chain} session object
comparable to \texttt{self.\_session} on the C++/pure-Python backends for
\texttt{kpmdmrg.py} to dispatch to directly), then hands the whole moment
loop to Julia in one call. Measured directly on a 30-site Heisenberg
chain: eliminating the per-step round trip only closed part of the gap
to the compiled ITensor v3 backend (v3: 23.4s, Julia native loop warm:
34.0s, vs.\ 37.7s for the earlier per-step Python loop) --- most of the
remaining time is genuine Julia-side tensor-contraction/truncation cost
(\texttt{alg="densitymatrix"} SVD-based \texttt{add}/\texttt{contract}),
not Python$\leftrightarrow$Julia marshaling overhead.

\textbf{Follow-up optimization}: \texttt{kpm\_moments\_full}/
\texttt{kpm\_moments\_accelerated} originally called
\texttt{mpsalgebra.jl}'s \texttt{applyoperator()} once per Chebyshev
step, which does \emph{two} truncated MPO-MPS contractions per call
(\texttt{contract(A,psi)} plus a second contraction against an identity
MPO --- the same "fixes a bug" workaround documented above for
\texttt{tdvp.jl}'s \texttt{apply\_clean}, which switched to
\texttt{ITensorMPS}'s own higher-level \texttt{apply()} for the same
reason, on the TDVP side, to fix a \emph{different} bug).
\texttt{mpsalgebra.jl}'s \texttt{apply\_op()} does the same
swap, cutting one of the two contractions per moment (not a clean
$2\times$, since the identity-MPO contraction is against a trivial
bond-dimension-1 operator, much cheaper than the contraction against the
real, bond-growing scaled Hamiltonian). Measured directly on the same
30-site chain: 34.5s $\to$ 28.4s warm ($\sim$18\% faster), narrowing the
gap to v3 from $\sim$1.47$\times$ to $\sim$1.21$\times$ slower. Peak
positions confirmed unchanged against the existing ED cross-check after
the change. \texttt{apply\_op()} is shared between \texttt{kpm.jl}'s
recursion and \texttt{tdvp.jl}'s \texttt{apply\_clean()} (which layers
its own \texttt{orthogonalize!()} on top, needed for
\texttt{tdvp()}/\texttt{dmrg()} inputs but not for KPM's) --- it lives in
\texttt{mpsalgebra.jl}, loaded before both, rather than as two
independently maintained near-duplicates. \texttt{mpsalgebra.jl}'s
\texttt{summps()} (the \texttt{add()} wrapper the Chebyshev recursion
sums consecutive Chebyshev vectors with) now also takes an explicit
\texttt{cutoff} argument: \texttt{add()}'s own default ($10^{-15}$) is
three orders tighter than dmrgpy's usual configured cutoffs, so every
moment step used to silently keep far more Schmidt values than
\texttt{self.kpmcutoff} called for, up to \texttt{kpmmaxm}. The KPM seed
vectors (\texttt{psi1}/\texttt{psi2} in
\texttt{mpsjulialive/dynamics.py}) are now also built via
\texttt{apply\_op()} rather than the older \texttt{applyoperator()},
closing the one spot the original optimization pass left on the old
primitive.

Real-time TDVP evolution (\texttt{timedependent.py}'s
\texttt{evolve\_and\_measure\_dmrg}/\texttt{evolution\_dmrg\_DC}, submode
\texttt{"TD"}) is implemented the same way ---
\texttt{mpsjulialive/tdvp.jl}'s \texttt{evolve\_and\_measure\_tdvp}/
\texttt{quench\_tdvp} run the whole nt-step trajectory in one Julia call,
driving \texttt{ITensorMPS.jl}'s own \texttt{tdvp()} (already exported by
the \texttt{ITensorMPS} dependency; no separate \texttt{ITensorTDVP.jl}
package needed). One real-time step is $\exp(-i \cdot dt \cdot H)$, i.e.\
\texttt{tdvp()}'s complex evolution parameter is \texttt{-im*dt}. Two
real, non-obvious \texttt{ITensorMPS}/\texttt{ITensors} bugs had to be
worked around to get this correct, both confirmed by direct inspection
of the intermediate tensors' index structure rather than guessed at from
the stack trace alone:

\begin{itemize}
\item \texttt{mpsalgebra.jl}'s \texttt{applyoperator()}
  (\texttt{contract(A,psi)} plus a second contraction against an
  identity MPO, "this fixes a bug" --- a \emph{Site}-index priming
  issue) leaves the \emph{Link} indices of its result carrying a stale
  nonzero prime level. This is invisible to KPM (which only ever feeds
  \texttt{applyoperator()}'s output back into more
  \texttt{applyoperator}/\texttt{summps}/\texttt{inner} calls), but
  \texttt{tdvp()}'s internal environment bookkeeping (\texttt{ProjMPO})
  cannot handle it: it aborts deep inside
  \texttt{KrylovKit.expintegrator} with "... but the indices are not
  permutations of each other" on the affected Link index.
  \texttt{tdvp.jl}'s \texttt{apply\_clean()} uses \texttt{ITensorMPS}'s
  higher-level \texttt{apply()} instead (\texttt{contract()} is the
  lower-level primitive both \texttt{applyoperator()} and
  \texttt{apply()} are built on, but only \texttt{apply()} avoids the
  stale prime).
\item Whenever \texttt{tdvp()}'s input MPS did not come straight from
  \texttt{dmrg()} (e.g.\ \texttt{evolution\_ABA()} first applies an
  operator via \texttt{mpsjulialive/mpo.py}'s \texttt{MPO.\_\_mul\_\_},
  which still goes through \texttt{applyoperator()} for other callers'
  sake), the same stale Link prime reappears --- and, confirmed
  directly, \texttt{orthogonalize!()} alone does \emph{not} clear it.
  \texttt{evolve\_and\_measure\_tdvp()} therefore explicitly
  \texttt{noprime(copy(wf),"Link")}s (never mutating the caller's own
  MPS) before \texttt{orthogonalize!()}-ing and starting the TDVP loop.
\end{itemize}

Both paths were validated directly against exact diagonalization (quench
from a N\'eel-favoring field to a Heisenberg chain, mirroring
\texttt{examples/time\_evolution/tdvp\_VS\_ED\_time\_evolution}):
\texttt{evolve\_and\_measure} agreed with ED to $\sim 8 \times 10^{-7}$
on a 4-site chain, and \texttt{evolution\_DC}'s quench correlator agreed
to $\sim 9 \times 10^{-11}$ on a 6-site chain.

Excited states (\texttt{get\_excited\_states}/\texttt{get\_excited},
$n>1$, Hermitian) and the single-site reduced density matrix
(\texttt{get\_rdm}) are implemented the same way:
\texttt{mpsjulialive/excited.jl}'s \texttt{excited\_states\_dmrg} runs
the whole $n$-state loop in one Julia call (one new random-start warm
state per additional excited state, deflated against every wavefunction
found so far via \texttt{ITensorMPS.jl}'s own orthogonality-penalty
\texttt{dmrg(H, wfs, psi0, sweeps; weight)}, mirroring
\texttt{mpscpp3/chain\_session.h}'s \texttt{Chain::excited\_states} and
\texttt{pyitensor/chain.py}'s \texttt{excited\_states});
\texttt{mpsjulialive/densitymatrix.jl}'s \texttt{reduced\_dm} is a
direct Julia port of \texttt{pyitensor/chain.py}'s \texttt{reduced\_dm}
(itself a port of \texttt{Chain::reduced\_dm}), including its "divide by
the norm squared, not its square root" quirk, preserved for
cross-backend consistency (in practice a no-op, since the ground state
is essentially always already unit-norm). Validated directly: excited
energies match the golden regression values in
\texttt{tests/test\_excited\_states.py} to $\sim 3 \times 10^{-15}$ on a
4-site chain, and \texttt{get\_rdm} matches the existing
backend-agnostic \texttt{reduced\_dm\_projective} to
$\sim 1 \times 10^{-15}$. $n=1$ and non-Hermitian excited states already
worked before this --- they route through \texttt{gs\_energy()}/the
generic MPS Arnoldi method (default \texttt{algebra/arpacktk.py}'s
IRAM, via \texttt{mpsalgebra.mps\_excited\_states}; \texttt{algebra/
arnolditk.py}'s \texttt{mpsarnoldi} remains available for comparison),
neither of which is backend-specific.

Bond entanglement entropy (\texttt{MPS.get\_bond\_entropy}, used e.g.\ by
\texttt{entropy.py}'s \texttt{central\_charge}) was missing entirely on
the Julia backend (\texttt{mpsjulialive/mps.py}'s \texttt{MPS} class had
no \texttt{get\_bond\_entropy}/\texttt{get\_site\_entropy} methods at
all --- a plain \texttt{AttributeError}, not a crash inside a dispatch
branch). \texttt{mpsjulialive/entropy.jl}'s \texttt{bond\_entropy}
computes it the standard MPS way, via SVD of the two-site tensor at that
bond (mirrors \texttt{pyitensor/chain.py}'s \texttt{bond\_entropy},
itself a port of \texttt{Chain::bond\_entropy}) --- much cheaper than
the generic \texttt{get\_correlation\_entropy}/\texttt{"explicit"}-dmmode
fallback the Julia backend already had (still the only option for
multi-site correlation entropy, just not for a single bond). Validated
directly: a 2-site Heisenberg singlet gives exactly $\ln 2$ (the
analytically known Bell-pair entanglement), and on a 6-site chain the
bond entropy at the first bond exactly matches the existing generic
\texttt{get\_site\_entropy} (which computes the same quantity a
completely different way, through \texttt{reduced\_dm\_projective}).

The CVM dynamical-correlator submode
(\texttt{get\_dynamical\_correlator(..., submode="CVM")},
\texttt{cvm.py}) needed almost no Julia-specific code at all: its
correction-vector CG solve (\texttt{cvm.py::cvm\_correction\_vector}) is
already implemented purely with backend-agnostic MPS/MPO algebra
(\texttt{self.toMPO()}, MPS \texttt{+}/\texttt{-}/scalar-\texttt{*},
\texttt{.dot()}), which already works against \texttt{julia\_live}'s
\texttt{mpsjulialive.mpo.MPO}/\texttt{mpsjulialive.mps.MPS} classes
without any changes. The only thing standing in the way was that
\texttt{cvm.py} unconditionally called
\texttt{self.\_session.set\_sweep\_params(...)} to point every MPO
application at \texttt{cvm\_maxm} rather than the DMRG \texttt{maxm} (no
\texttt{\_session} object exists for \texttt{julia\_live} to call this
on); \texttt{cvm.py::\_set\_cvm\_sweep\_params} now temporarily
overrides \texttt{self.maxm} instead for that backend, since that's what
\texttt{mpsjulialive/mpo.py}/\texttt{mps.py} actually read. Validated
against ED with the same tolerance
\texttt{examples/dynamical\_correlator/dynamical\_correlator\_VS\_ED/main.py} uses for the
C++/pure-Python backends ($10^{-6}$) --- matched to
$\sim 2 \times 10^{-15}$ here, in 3 CG iterations per frequency point.

Wiring this up surfaced a real bug in
\texttt{mpsjulialive/dynamics.py::get\_dynamical\_correlator}: its first
version declared \texttt{name=}/\texttt{delta=}/\texttt{es=} as its own
named parameters (with KPM-flavored defaults) purely so it could compute
the KPM moments directly in the same function, which silently captured a
caller's \texttt{es=}/\texttt{name=}/\texttt{delta=} kwargs out of
\texttt{**kwargs} before they could reach
\texttt{cvm.dynamical\_correlator} --- confirmed directly,
\texttt{submode="CVM"} ran on a completely different, wrong frequency
grid as a result (CVM's own \texttt{np.linspace(0.,10.0,100)} default
instead of the caller's requested window). Fixed by splitting the KPM
implementation out into its own \texttt{\_kpm\_dynamical\_correlator}
and making the public
\texttt{get\_dynamical\_correlator(self,submode="KPM",**kwargs)} take no
named parameters of its own --- mirroring the top-level
\texttt{dynamics.py::get\_dynamical\_correlator}'s own signature, which
was never vulnerable to this because it never declared submode-specific
parameters either.

The TDZ dynamical-correlator submode (complex-time evolution + perturbative
real-axis reconstruction, \texttt{tdz.py}, arXiv:2311.10909) needed only
one new backend primitive, same as CVM: the whole algorithm is otherwise
generic MPS/MPO algebra (\texttt{self.toMPO()}, \texttt{.dot()}) already
working on \texttt{julia\_live}. That one primitive is a single
complex-time-step propagator
(\texttt{tdz.py::\_advance\_complex\_time\_step}, formerly gated to
\texttt{itensor\_version in (3,"python")}) ---
\texttt{mpsjulialive/tdvp.jl}'s \texttt{tdvp\_step} already generalizes
to it with \emph{no changes}, since \texttt{-im*dz} is exactly the right
formula whether \texttt{dz} is real (the
\texttt{evolve\_and\_measure\_tdvp}/\texttt{quench\_tdvp} case) or
genuinely complex (TDZ's per-step contour increment); only a thin Python
wrapper
(\texttt{mpsjulialive/timedependent.py::advance\_complex\_time\_step})
was needed.

Wiring TDZ up caught a second occurrence of the same stale-Link-prime bug
documented above for TDVP: \texttt{tdz.py}'s first evolved state,
\texttt{self.toMPO(A)*wf\_g}, goes through the same
\texttt{applyoperator()} path \texttt{evolution\_ABA()} does, and
\texttt{tdvp\_step} had only been sanitized against that inside
\texttt{evolve\_and\_measure\_tdvp}'s own wrapper loop, not inside
\texttt{tdvp\_step} itself --- so any \emph{other} direct caller (TDZ's
driver loop in \texttt{\_complex\_time\_correlator}) hit the identical
failure. Fixed by moving the \texttt{noprime(copy(psi),"Link")} +
\texttt{orthogonalize!()} sanitization into \texttt{tdvp\_step} itself,
so it runs on every step regardless of caller, instead of relying on
each call site to remember to pre-clean its input --- a small, one-time
cost next to the sweep itself. Validated against the same golden test
\texttt{tests/test\_dynamical\_correlator.py} uses for the other
backends: peak within $0.03$ of the exact 4-site Heisenberg gap
($0.658919$) --- landed at $0.68$.

The last two dynamical-correlator submodes needed no new Julia code at
all, only dispatch routing: \texttt{dcex.py} (submode \texttt{"EX"},
correlator via exact diagonalization in the excited-state subspace) only
calls \texttt{self.get\_excited\_states()} (already backend-agnostic,
see above) and generic \texttt{MultiOperator}/MPS algebra
(\texttt{.dot()}, \texttt{A*wf}) before dropping into plain NumPy/SciPy
(\texttt{eigh}); \texttt{distribution.py}'s maxent path (submode
\texttt{"maxent"}) is built the same way on top of
\texttt{vev.py::power\_vev}. Validated \texttt{"EX"} the same way
\texttt{tests/test\_dynamical\_correlator.py} already does for the other
backends --- cross-checked directly against \texttt{"KPM"} and
\texttt{"CVM"} on the same chain/operator/frequency grid, landing on the
exact same peak (diff $0.0$) rather than just the analytic gap.
\texttt{"maxent"} is wired up for parity but not actually functional on
\emph{any} backend right now --- \texttt{distribution.py::get\_distribution\_maxent}
imports \texttt{dmrgpy.maxenttk.pymaxent}, which doesn't exist in this
checkout (confirmed directly: \texttt{ModuleNotFoundError}), so it hits
the same \texttt{except: print("Not functional yet"); exit()} regardless
of \texttt{itensor\_version} --- a pre-existing, cross-backend gap, not
something introduced or left broken by this work.

The fermionic 4-point correlator tensor
(\texttt{MPS.get\_four\_correlation\_tensor}, $\langle
C^\dagger_i C_j C^\dagger_k C_l \rangle$) was missing the same way bond
entropy was --- no \texttt{get\_four\_correlation\_tensor} method at all
on \texttt{mpsjulialive/mps.py}'s \texttt{MPS} class --- and needed the
same fix: just add the method, delegating to
\texttt{entropytk/correlationentropy.py}'s default
\texttt{ctmode="explicit"} (a Python loop of \texttt{MultiOperator}
products + \texttt{.aMb()}/\texttt{.dot()}, already generic;
\texttt{ctmode="full"}, a native per-element AutoMPO build, mirrors
\texttt{pyitensor/chain.py}'s own \texttt{four\_correlation\_tensor} but
isn't ported to Julia). Validated directly against ED on the same
well-gapped free-fermion model \texttt{tests/test\_four\_point\_correlator.py}
uses (agreement to $\sim 2 \times 10^{-15}$).

\texttt{get\_four\_correlation\_tensor\_explicit}'s Python loop
(\texttt{entropytk/correlationentropy.py}) did not exploit the tensor's
own Hermitian symmetry ($\langle C^\dagger_i C_j C^\dagger_k
C_l\rangle^\dagger = C^\dagger_l C_k C^\dagger_j C_i$, so
\texttt{ct[i,j,k,l]} and \texttt{ct[l,k,j,i]} are complex conjugates)
the way \texttt{ctmode="full"}'s C++/pyitensor implementations already
did via their own \texttt{accelerate} flag --- added the same
\texttt{accelerate=True} default there too (skip one representative of
each \texttt{(current,conjugate)} pair, fill in its mirror from the
other), an exact, not approximate, speedup. This mattered for
\texttt{fermionchain.Spinful\_Fermionic\_Chain\_Native} (added
alongside the interleaved \texttt{Spinful\_Fermionic\_Chain}, see the
model table above): its \texttt{self.C}/\texttt{self.Cdag} are flat,
single-flavor-per-entry lists (\texttt{Cup}/\texttt{Cdn} interleaved as
mode $2i$/$2i+1$, matching \texttt{Spinful\_Fermionic\_Chain}'s own
indexing exactly, so the two classes' tensors compare index for index)
added purely so this already-generic \texttt{ctmode="explicit"} path
works unchanged for it.

\texttt{ctmode="full"} (the C++-accelerated path) was then added for
this class too, via a dedicated
\texttt{Chain::four\_correlation\_tensor\_spinful()}
(\texttt{mpscpp3/chain\_session.h}/\texttt{bindings.cc}) --- the
existing \texttt{Chain::four\_correlation\_tensor()} could not be
reused as-is, since it hardcodes the literal \texttt{"Cdag"}/\texttt{"C"}
operator names, undefined on ITensor's \texttt{ElectronSite} (only
\texttt{Cup}/\texttt{Cdn}/\texttt{Cdagup}/\texttt{Cdagdn} are). The new
method instead hands ITensor's own \texttt{AutoMPO} the flavor-resolved
names directly and relies on its built-in automatic fermionic-sign
insertion (\texttt{autompo.cc}'s
\texttt{isFermionic()}/\texttt{fermionicTerm()}, triggered by any
operator name starting with \texttt{'C'}) to do the Jordan-Wigner
threading --- a deliberate, one-off exception to this codebase's usual
rule of threading Jordan-Wigner strings explicitly at the Python level
for backend-agnosticism (see
\texttt{multioperatortk/jordanwigner\_spinful.py}); safe here only
because this calculation always builds and discards its own fresh,
self-contained \texttt{AutoMPO}, not a change to the general
Hamiltonian/MPO pipeline.
\texttt{entropytk/correlationentropy.py::get\_four\_correlation\_tensor\_cpp}
dispatches to whichever C++ method matches \texttt{type(wf.MBO)}.

Measured ($n=3,4,5,6,12$ orbitals):
\texttt{Spinful\_Fermionic\_Chain\_Native}'s \texttt{ctmode="full"} is
the fastest of all four combinations (native/interleaved $\times$
explicit/full) tried at every size, including $n=12$ (24 flat modes:
$\sim620$s vs $\sim890$s for \texttt{Spinful\_Fermionic\_Chain}'s own
\texttt{ctmode="full"}, a $\sim30\%$ win --- see
\texttt{examples/staticcorrelators/four\_correlation\_tensor\_spinful\_native}). Before
this C++-accelerated path existed, native's only option
(\texttt{ctmode="explicit"}) still beat interleaved's
\texttt{ctmode="full"} at $n=3..6$, by a margin that grew with $n$
there but did not continue to $n=12$ (the two
\texttt{ctmode="explicit"} numbers alone came back to essentially
tied, $\sim700$s each) --- adding \texttt{ctmode="full"} for this class
is what restores and extends the win at $n=12$. This remains the one
calculation checked so far where the native-site class wins at all,
since it is a static overlap computation rather than an iterative
two-site search, so it never pays the two-site combined-local-
dimension penalty documented in
\texttt{Spinful\_Fermionic\_Chain\_Native}'s own class docstring.

A third implementation, \texttt{ctmode="sweep"}
(\texttt{Chain::four\_correlation\_tensor\_sweep},
\texttt{mpscpp3/chain\_session.h}/\texttt{bindings.cc}, plain
non-native-spinful sites only), replaces \texttt{ctmode="full"}'s
$O(N^4)$ \emph{independent} AutoMPO builds with a single-sweep,
environment-reuse algorithm following the idea of
\href{https://github.com/ITensor/ITensorCorrelators.jl}{ITensorCorrelators.jl}.
Only the $N(N-1)(N-2)(N-3)$ pairwise-distinct-index entries go through
the fast sweep (nested loops over four strictly increasing sites
$a<b<c<d$, each level extending a running left environment by exactly
one site rather than rebuilding it, with both \texttt{Cdag}/\texttt{C}
tried at each of the four operator sites); the remaining, subdominant
$O(N^3)$ repeated-index entries fall back to the same per-tuple AutoMPO
method \texttt{ctmode="full"} uses, since those need same-site
multi-operator products (e.g.\ $C^\dagger_i C_i = n_i$) the sweep is not
built to produce. Reordering the four fermionic operators from the
abstract $(i,j,k,l)$ order into physical site order picks up the
standard fermion-anticommutation sign (parity of the reordering
permutation) \emph{plus} one further correction: of the six possible
\texttt{Cdag}/\texttt{C} type patterns across the four sorted sites, the
two that strictly alternate (\texttt{Cdag,C,Cdag,C} or
\texttt{C,Cdag,C,Cdag}) need no extra sign, but the other four (which
have at least one pair of \emph{adjacent same-type} operators) need one
more flip on top of the plain permutation parity --- this was found
empirically (every entry with such a pattern disagreed with
\texttt{ctmode="full"} by exactly this sign, for every site quadruple
and chain length tried, before the correction was added) rather than
purely re-derived by hand from Jordan-Wigner strings, though it traces
to the same extra $-1$ that \texttt{jordanwigner.py}'s
\texttt{CC()}/\texttt{CdagCdag()} carry but
\texttt{CdagC()}/\texttt{CCdag()} don't (reordering two \emph{same-type}
operators into JW-canonical form costs one more local $F/A^\dagger$ or
$F/A$ anticommutation than a mixed pair does). Validated to machine
precision against \texttt{ctmode="full"} and to ED solver tolerance
across $n=4..8$ (every index tuple, not just a sample) and spot-checked
at $n=16$; measurably faster than \texttt{ctmode="full"} at every size
tried, by a margin growing with $n$ ($\sim1.4\times$ at $n=6$ up to over
$2.5\times$ at $n=16$ for a nearest-neighbor Hubbard chain at fixed
\texttt{maxm=60} --- absolute numbers are noisy across runs on a shared
machine (measured $2.8$--$3.4\times$ at $n=16$ across separate runs),
but the \emph{trend}, measured within any single run, consistently
holds --- see
\texttt{examples/staticcorrelators/four\_correlation\_tensor\_sweep\_VS\_full},
whose Part 3 loop actually runs and reproduces this trend, rather than
reporting only development-time measurements that would not be
reproducible from the checked-in tree; $n=20$ was also checked
informally during development, $\sim4.4\times$, but is not part of the
shipped example since a single $n=20$ \texttt{ctmode="full"} call alone
takes $\sim5$ minutes). \texttt{accelerate} (default \texttt{True}) only
gates the subdominant repeated-index fallback here --- unlike
\texttt{ctmode="full"}'s own \texttt{accelerate}, which skips $\sim$half
of the \emph{dominant} per-tuple AutoMPO builds via conjugate-pair
symmetry, this method's dominant pairwise-distinct sweep pays its cost
(the shared environment sweep, and six cheap per-pattern
\texttt{eltC()} evaluations) once regardless of how many output entries
a given leaf value gets scattered into, so there is no equivalent saving
to skip there --- confirmed directly (\texttt{accelerate=True} vs.\
\texttt{False} gives identical wall-clock on the dominant part).

\texttt{four\_correlation\_tensor\_sweep()} originally renormalized its
internal copy of \texttt{wf} before the fast-sweep computation
(\texttt{psi /= innerC(psi,psi).real()}, mirroring \texttt{reduced\_dm()}'s
own convention), but the repeated-index fallback loop calls
\texttt{innerC(wf,...)} on the raw, un-renormalized \texttt{wf} --- a
real bug, caught by code review, not by any of the (unit-norm-only)
validation above: for any non-unit-norm input MPS (e.g.\ \texttt{wf*c}
for a scalar $c$, as opposed to an already-converged DMRG ground state,
which is unit-norm to solver precision and so did not expose it), the
two halves of one output tensor came back scaled inconsistently with
each other \emph{and} with \texttt{ctmode="full"}/\texttt{"explicit"}.
Fixed by dropping the renormalization entirely, matching
\texttt{four\_correlation\_tensor()}'s own convention (no enforced
normalization, caller's responsibility) exactly --- confirmed with a
direct repro (\texttt{wf*3.0}: \texttt{ctmode="full"} scales every entry
by $3^2=9$ as expected; \texttt{ctmode="sweep"} did too only on the
repeated-index entries, and by $1/9$ on the pairwise-distinct ones,
before the fix). The same bug, same fix, applies to the
\texttt{pyitensor/chain.py} port below.

Not ported to \texttt{Spinful\_Fermionic\_Chain\_Native}: that class's
\texttt{ctmode="full"} gets its Jordan-Wigner threading for free from
ITensor's own \texttt{AutoMPO} acting on flavor-resolved operator names,
since two flavors share one physical site there --- a different
threading problem than this sweep's per-flat-mode gap rule solves, not
attempted here.

The same algorithm was then ported to \texttt{pyitensor/chain.py}'s
\texttt{Chain.four\_correlation\_tensor\_sweep} for
\texttt{itensor\_version="python"} --- close to a mechanical
transcription of the C++ version, including
\texttt{\_four\_pt\_perm\_table()} (the same sign/index-role table,
generated with \texttt{itertools.permutations} instead of a hand-rolled
\texttt{next\_permutation}), with one adaptation: this engine never
primes \texttt{Link}-tagged indices to keep a self-overlap's bra and ket
from colliding the way \texttt{chain\_session.h}'s
\texttt{dag(prime(psi.A(site), TagSet("Link")))} does (see
\texttt{mpsalgebra.py}'s \texttt{inner()}) --- instead the bra side is
built once per outer site $a$ via \texttt{mpsalgebra.py}'s own
\texttt{\_fresh\_link\_copy(psi)} (freshly relabeled \texttt{Link}
indices, physical indices untouched), and the running environment
carries one leg from the \emph{ket}'s own unrelabeled links and one from
this fresh bra copy's links; the boundary \texttt{delta} ``shortcut''
tensors are built from \texttt{commonIndex} on each side (ket vs.\
fresh-bra) rather than from a single primed/unprimed pair. Matched the
C++ version and ED to machine precision immediately (0 mismatches out of
every distinct- and repeated-index tuple tried, $n=5$), confirming the
ported sign table is correct and the fresh-link-copy substitution is a
faithful translation, not just an analogous-looking one. Real but
smaller speedup than the C++ version ($\sim1.2$--$1.4\times$ at
$n=5..7$, vs.\ $\sim1.4$--$4.4\times$ at $n=6..20$ for
\texttt{itensor\_version=3}): pure-Python loop/function-call overhead is
a bigger fraction of the per-step cost here than the underlying NumPy
linear algebra, so there is less for environment reuse to save on
relative to the total.

\texttt{get\_four\_correlation\_tensor(ctmode=...)}'s default changed
from a fixed \texttt{ctmode="explicit"} to \texttt{ctmode=None}, resolved
per-wavefunction by the new
\texttt{\_four\_correlation\_tensor\_default\_ctmode()}
(\texttt{entropytk/correlationentropy.py}): \texttt{"sweep"} whenever it
applies (\texttt{itensor\_version} in \texttt{(3,"python")},
non-native-spinful), else \texttt{"full"} whenever it applies
(\texttt{itensor\_version} in \texttt{(2,3,"python")}, or native-spinful
under \texttt{itensor\_version=3}), else \texttt{"explicit"}. An explicit
\texttt{ctmode=} argument is unaffected --- it is still a hard request
that raises if unavailable, not a hint the resolver can override.
ED-backed wavefunctions (\texttt{edtk/edchain.py},
\texttt{pyfermion/mbfermion.py}) never reach the resolver at all: they
hardcode \texttt{ctmode="explicit"} themselves, as before. Verified the
resolver picks the right mode across every backend actually reachable
from the public API (\texttt{itensor\_version} \texttt{2} falling back
to ED when the C++ extension is not compiled, \texttt{3},
\texttt{"python"}, and \texttt{Spinful\_Fermionic\_Chain\_Native}), and
made it use \texttt{getattr()}/\texttt{hasattr()} defensively rather
than assume \texttt{wf.MBO} always carries
\texttt{.itensor\_version}/\texttt{.\_session} --- cheap insurance since
the function is easy to reach with an unusual \texttt{wf.MBO} from
outside the two call sites it is actually designed for.

Investigated whether \texttt{get\_correlation\_matrix}'s default
\texttt{dmmode="fast"}
(\texttt{entropytk/correlationentropy.py::correlation\_matrix\_fast}, $n$
MPO applications total rather than \texttt{explicit}'s $n(n+1)/2$
two-operator products) could also work on \texttt{julia\_live} ---
\texttt{mpsjulialive/mps.py}'s \texttt{get\_correlation\_entropy}/
\texttt{get\_correlation\_entropy\_density} already force
\texttt{dmmode="explicit"}, and this turned out to be necessary, not an
arbitrary choice: \texttt{dmmode="fast"} builds an MPO for a single bare
fermionic operator (\texttt{Cdag[i]}, odd particle-number parity) before
combining it with another, and \texttt{ITensorMPS.jl}'s OpSum-to-MPO
compiler rejects that outright ("Parity-odd fermionic terms not yet
supported by OpSum to MPO conversion", confirmed directly) ---
\texttt{explicit} avoids the problem by always building the two-operator
\emph{product} ($C^\dagger_i C_j$, even parity) before ever converting
to an MPO. Not pursued further: it's a real \texttt{ITensorMPS.jl}
limitation, not a bug on this side, and \texttt{explicit} is already
fast in absolute terms (0.29s for a 10-site chain's full correlation
matrix) --- a from-scratch fermionic-JW-string MPO builder to work
around it would be a disproportionate amount of new code for a
performance-only win on an already-cheap operation.

A final sweep of every remaining \texttt{self.\_session} call site
reachable from \texttt{julia\_live} (grepped across the package, checked
against what each call site's own top-level dispatcher actually routes
for that \texttt{itensor\_version}) turned up two more real gaps,
deliberately left unfixed rather than pursued further, since both fall
outside the "already-generic code, just needs a dispatch branch or a
small native Julia primitive" pattern every fix above followed:

\begin{itemize}
\item \textbf{\texttt{vev.py}'s \texttt{npow=} kwarg}
  (\texttt{sc.vev(op, npow=2)}, computes $\langle X^2\rangle$/$\langle
  X^n\rangle$ via repeated MPO application instead of building the
  $O(n^2)$-term squared operator directly) --- \texttt{mpsjulialive/vev.py}'s
  \texttt{vev(MBO,MO)} takes no \texttt{**kwargs} at all, so this raises
  a plain \texttt{TypeError} for \texttt{julia\_live} rather than
  silently misbehaving. The one example exercising it,
  \texttt{examples/utilities/power\_vev/main.py}, sets
  \texttt{sc.itensor\_version = "julia"} directly (bypassing
  \texttt{setup\_julia()}/\texttt{\_reset\_dmrg\_state()} entirely) ---
  the legacy, already-inert subprocess-based backend documented below,
  not \texttt{julia\_live} --- so it wouldn't validate this fix even if
  made.
\item \textbf{\texttt{get\_distribution}/\texttt{kpmdmrg.py::general\_kpm}}
  (spectral distribution of an arbitrary operator, not just the
  Hamiltonian; has example coverage in four
  \texttt{examples/magnetization/magnetization\_distribution*} scripts and
  \texttt{examples/dynamical\_correlator/dynamical\_correlator\_shift},
  but no \texttt{pytest} coverage) --- unlike everything above, this is
  not a wire-up-already-generic-code fix:
  \texttt{kpmdmrg.py::general\_kpm\_moments} depends on
  \texttt{scale\_operator()}, which calls
  \texttt{self.lowest\_eigenvalue()}/\texttt{self.bandwidth()}
  (\texttt{manybodychain.py}), both of which route through
  \texttt{self.clone()} $\to$ \texttt{deepcopy(self)} --- and
  \texttt{julia\_live}'s live Julia session handles
  (\texttt{self.jlsites}, every \texttt{MPS.jlmps}/\texttt{MPO.jlmpo})
  are not established to be deepcopy-safe (this is exactly why
  \texttt{dynamics.py::\_max\_energy\_bound}/\texttt{excited.py}'s
  energy-bound helpers were written as their own mutate-and-restore
  functions instead of calling \texttt{bandwidth()}/
  \texttt{lowest\_eigenvalue()} directly, see their docstrings). Doing
  this properly means generalizing that same mutate-and-restore pattern
  from "the Hamiltonian, negated" to "an arbitrary caller-supplied
  operator" --- real, but new design work rather than a small addition,
  so left as a documented follow-up rather than implemented here.
\end{itemize}

\textbf{Fixes from an independent multi-angle code review of the whole
Julia backend} (8 finder passes + per-finding verification, run once the
above was otherwise feature-complete): six real, confirmed issues, all
fixed directly rather than left as follow-ups ---

\begin{itemize}
\item \texttt{cvm.py}'s \texttt{\_set\_cvm\_sweep\_params} (the
  \texttt{julia\_live} equivalent of
  \texttt{self.\_session.set\_sweep\_params}) is now a context manager
  (\texttt{\_cvm\_sweep\_params}) instead of a plain
  set-and-return-the-old-value function: \texttt{self.maxm} is restored
  via \texttt{finally}, so a CG blow-up mid-loop (or a standalone call to
  the public \texttt{cvm\_correction\_vector}, which previously never
  restored it at all) can no longer leave \texttt{self.maxm} permanently
  stuck at \texttt{cvm\_maxm}.
\item \texttt{dynamics.py}'s \texttt{julia\_live} branch now checks
  \texttt{self.is\_hermitian(...)} before dispatching, the same way the
  \texttt{(2,3,"python")} branch already does --- previously a
  non-Hermitian Hamiltonian silently ran the Hermitian-only KPM/CVM/TDZ
  math and returned numerically wrong output. There's no non-Hermitian
  route to fall back to for \texttt{julia\_live} yet (that needs
  \texttt{applyinverse()}, C++-only today), so this raises a clear
  \texttt{NotImplementedError} instead.
\item \texttt{mpsjulialive/excited.py}'s \texttt{get\_excited\_states\_dmrg}
  now honors \texttt{self.excited\_gram\_schmidt} (previously silently
  dropped --- \texttt{excited.jl} has no Gram-Schmidt step of its own),
  applying the same generic, backend-agnostic
  \texttt{algebra.arnolditk.gram\_smith} the top-level
  \texttt{get\_excited\_states}'s own \texttt{purify=True} path already
  uses, and recomputing energies against the orthogonalized states to
  match.
\item \texttt{mpsjulialive/dynamics.py}'s \texttt{\_max\_energy\_bound}
  now restores \texttt{self.maxm}/\texttt{self.nsweeps}/
  \texttt{self.hamiltonian} in a \texttt{finally} block --- previously a
  \texttt{self.gs\_energy()} failure partway through (a real possibility
  for a live juliacall session) would leave the chain object with a
  negated Hamiltonian and clamped bond dimension/sweep count for every
  later call.
\item \texttt{mpsjulialive/tdvp.jl}'s \texttt{evolve\_and\_measure\_tdvp}
  now explicitly renormalizes every step (matching \texttt{quench\_tdvp}'s
  own existing pattern, and \texttt{mpscpp3/chain\_session.h}'s
  \texttt{tdvp\_step}, which passes \texttt{"DoNormalize",true}) ---
  \texttt{tdvp\_step}'s own \texttt{normalize=false} stays unchanged,
  since \texttt{tdz.py}'s \texttt{advance\_complex\_time\_step} also
  calls \texttt{tdvp\_step} directly and its physically-meaningful norm
  decay (the whole point of TDZ's damping mechanism, arXiv:2311.10909)
  must not be renormalized away.
\item \texttt{mpsjulialive/tdvp.jl}'s \texttt{tdvp\_step} no longer forces
  \texttt{mindim=2}: that floor was a guessed fix for a bug later found
  to be the stale-Link-prime issue (fixed separately, see above) and was
  never confirmed necessary: neither \texttt{mpscpp3}'s own
  \texttt{tdvp\_step} nor \texttt{pyitensor/tdvp.py}'s \texttt{svd} calls
  force a floor above ITensor's own default of 1, and forcing 2 could pad
  in a spurious singular value at a bond whose true Schmidt rank
  genuinely is 1 (e.g.\ a near-product state on a small chain).
\end{itemize}

One more finding was fixed as a cleanup rather than a correctness bug:
\texttt{mpsjulialive/excited.jl}'s \texttt{excited\_state\_dmrg} had
copy-pasted the \texttt{Sweeps}-construction + stdout-silencing
boilerplate a third time (\texttt{get\_gs.jl} already had two
near-identical copies); both are now built on shared
\texttt{make\_sweeps}/\texttt{run\_quiet} helpers in \texttt{get\_gs.jl}.

A seventh candidate (\texttt{densitymatrix.jl}/\texttt{entropy.jl}
lacking a bounds check when \texttt{site}/\texttt{b} is the chain's last
site) was investigated and found to be a pre-existing,
deliberately-preserved limitation shared identically by
\texttt{pyitensor/chain.py} and \texttt{mpscpp3/chain\_session.h} (see
\texttt{reduced\_dm}'s own docstring above) --- not a new risk, so left
as-is.

(A separate, older subprocess-based Julia path,
\texttt{itensor\_version="julia"} via \texttt{juliarun.py}, is not
reachable through the normal public API and should be treated as
legacy/inert.)

\textbf{A second review round} (run against the KPM optimization + the
six fixes above together) found two of those fixes were themselves
incorrect, plus several smaller cleanup items --- all fixed, and this
time backed by persisted regression coverage
(\texttt{tests/test\_julia\_live.py}, \texttt{julia\_live}'s first,
since the first round's validation was ad hoc and non-committed):

\begin{itemize}
\item \texttt{evolve\_and\_measure\_tdvp}'s new renormalization (previous
  round, above) forced the state to \emph{unit} norm every step instead
  of preserving its own starting norm --- diverging from
  \texttt{quench\_tdvp}'s existing \texttt{norm0}-target pattern in the
  same file, and silently rescaling every result by $1/\lVert wf\rVert^2$
  whenever the input wasn't already unit-norm (e.g.\
  \texttt{evolution\_ABA()}'s \texttt{wfA = A*wf}, for any non-unitary
  \texttt{A}). Now captures \texttt{norm0 = }$\lVert wf\rVert$ once at
  entry and renormalizes to that every step, matching
  \texttt{quench\_tdvp} and the pure-Python reference's physical
  behavior.
\item The non-Hermitian guard added to \texttt{dynamics.py}'s
  \texttt{julia\_live} branch (previous round, above) ran \emph{before}
  submode dispatch, so it unconditionally blocked every submode ---
  including \texttt{submode="EX"}, which already has a working,
  itensor\_version-agnostic non-Hermitian path (\texttt{dcex.py}
  $\to$ \texttt{excited.py}'s \texttt{excited\_states\_non\_hermitian},
  not gated on \texttt{itensor\_version} at all) and
  \texttt{submode="maxent"}. The guard now only fires for
  \texttt{submode in ("KPM","CVM","TDZ")}, matching what its own error
  message already claimed.
\item Both renormalization sites now use \texttt{norm(psi)} instead of
  \texttt{sqrt(abs(inner(psi,psi)))}: \texttt{tdvp\_step()}'s own
  \texttt{orthogonalize!(psi,1)} guarantees a well-defined orthogonality
  center on its output, so \texttt{ITensorMPS}'s \texttt{norm()}
  short-circuits to a cheap $O(D^2)$ single-tensor norm instead of a full
  $O(N{\cdot}D^3)$ chain contraction --- confirmed via a live benchmark
  ($N=40$, $D=60$: $\sim$4ms vs $\sim$1.18s per call).
\item \texttt{mpsjulialive/dynamics.py}'s \texttt{\_same\_mps} had an
  unused \texttt{jlsites} parameter left over from the same cleanup that
  removed it from its sibling \texttt{\_kpm\_moments\_full}/
  \texttt{\_kpm\_moments\_accelerated} --- dropped, and
  \texttt{same\_mps}/\texttt{summps} on the Julia side gained the
  \texttt{cutoff} argument described above.
\end{itemize}

The commit that introduced this round's fixes claimed "new targeted
tests" without actually adding any (\texttt{tests/}/\texttt{examples/}
were untouched) --- \texttt{tests/test\_julia\_live.py} now covers
exactly the four behaviors that commit described validating: CVM's
\texttt{maxm} restoration, the non-Hermitian dispatch split,
\texttt{evolution\_ABA}'s norm preservation, and the KPM peak position
after the \texttt{apply\_op}/\texttt{summps} changes. The whole file is
skipped if \texttt{juliacall}/Julia isn't available, the same way
\texttt{tests/} already skips itensor\_version 2/3 when the corresponding
compiled extension isn't present.

\subsection{Supporting \texttt{*tk} packages}

Functionality is generally split into a top-level module (the public
method) and a \texttt{<topic>tk/} package with the implementation, e.g.\
\texttt{correlator.py} / \texttt{fermionchaintk/staticcorrelator.py},
\texttt{dynamics.py} / \texttt{dynamicstk/}, \texttt{entropy.py} /
\texttt{entropytk/}, \texttt{mpsalgebra.py} / \texttt{mpsalgebratk/}.
When changing behavior for a specific feature, check both the thin
dispatch module and its \texttt{tk} counterpart.

\subsection{KPM / dynamical correlators}

Dynamical correlators and generic operator distributions are computed
with the Kernel Polynomial Method (\texttt{kpmdmrg.py};
\texttt{Chain::kpm\_dynamical\_correlator} / \texttt{Chain::general\_kpm}
in \texttt{chain\_session.h} on the C++ side;
\texttt{algebra/kpm.py}-style moment recursion on the ED side)
rather than by direct spectral decomposition, since exact
diagonalization of the full spectrum is infeasible for large chains.

\subsection{Non-Hermitian KPM (NH-KPM)}

For a non-Hermitian Hamiltonian, \texttt{dynamics.py}/
\texttt{edtk/dynamics.py}'s \texttt{submode="KPM"} gate routes to
\texttt{nonhermitian/kpm.py} instead of \texttt{kpmdmrg.py}/
\texttt{algebra/kpm.py}'s Hermitian moment machinery (previously every
submode fell through unconditionally to the correction-vector fallback
\texttt{nonhermitian.dynamics.dynamical\_correlator\_non\_hermitian},
which assumes $A^\dagger=B$ and never used the left ground state at
all). NH-KPM is a port of the biorthogonal Chebyshev-moment algorithm in
NHKPM.jl\footnote{\url{https://github.com/GUANGZECHEN/NHKPM.jl}}
(Phys.\ Rev.\ Lett.\ \textbf{130}, 100401): rather than a single
self-dual ground state, it uses the biorthogonal pair
\texttt{(psi\_L,psi\_R)} already produced by NH-DMRG (\S4.4,
\texttt{nhdmrg()}/\texttt{gs\_energy\_nhdmrg}) on the DMRG side, or
\texttt{algebra.biorthogonal\_ground\_state} (a small addition to
\texttt{algebra.py} diagonalizing both \texttt{h} and
\texttt{dagger(h)} and matching/normalizing the pair so that
$\langle\psi_L|\psi_R\rangle=1$) on the ED side.

The algorithmic difference from Hermitian KPM: a non-Hermitian rescaled
operator $\tilde H(z)=(z\,\mathrm{Id}-H)/E_{\max}$ has no real spectrum,
so the ordinary single-operator Chebyshev recursion (valid once $H$ is
rescaled onto $[-1,1]$) does not apply. NH-KPM instead builds a
\emph{coupled} forward/adjoint recursion from both $\tilde H(z)$ and
$\tilde H(z)^\dagger$ (\texttt{get\_mu\_n\_nh}/
\texttt{spec\_from\_moments\_nh} in \texttt{algebra/kpm.py},
\texttt{Chain::nhkpm\_moments} in \texttt{mpscpp3/chain\_session.h},
ported line-for-line from the reference's
\texttt{get\_vn\_NH}/\texttt{get\_mu\_n\_NH}/\texttt{get\_spec\_kpm\_NH}).
The practical consequence: unlike Hermitian KPM (moments computed once,
reused at every frequency via cheap Chebyshev polynomial evaluation),
NH-KPM's expansion operator depends on $z$ itself, so
\texttt{nonhermitian/kpm.py} rebuilds the MPO/matrix and recomputes the
full moment recursion at \emph{every} requested frequency point ---
mirroring the reference algorithm's own cost profile, not a missed
optimization. \texttt{E\_max} must be supplied by the caller: there is
no automatic spectral-radius estimator yet for a non-Hermitian operator
(\texttt{chain\_session.h}'s \texttt{maximum\_energy()}, used for the
Hermitian case, runs ordinary variational DMRG on $-H$ and is
meaningless once $H$ is not Hermitian).

On the DMRG side, the $(z\,\mathrm{Id}-H)/E_{\max}$ operator and its
dagger are built once per frequency as ordinary \texttt{MultiOperator}s
in Python (\texttt{z*identity() - self.hamiltonian}, then
\texttt{.get\_dagger()}), converted via the existing
\texttt{to\_terms()}/\texttt{build\_mpo} machinery --- no new MPO
arithmetic was needed in C++, only the new coupled-recursion primitive
\texttt{Chain::nhkpm\_moments} (public, alongside \texttt{general\_kpm})
and its \texttt{bindings.cc} wrapper. One non-obvious bookkeeping fix was
needed to reuse NH-DMRG's own state: \texttt{gs\_energy\_nhdmrg}
renormalizes \texttt{wf0=psir.normalize()} to unit norm but leaves
\texttt{nh\_left\_wf} at NH-DMRG's own $\langle\psi_L|\psi_R\rangle=1$
scale, so the two attributes no longer satisfy that biorthogonal
relation together; \texttt{nonhermitian.kpm.dynamical\_correlator\_nhkpm}
restores it explicitly
(\texttt{psil = psil/conj(psil.dot(psir))}) before using the pair,
rather than assuming the convention still holds after
\texttt{gs\_energy()}.

Implemented so far for the ED backend, \texttt{itensor\_version=3}, and
\texttt{itensor\_version="python"} (\texttt{itensor\_version=2} raises
\texttt{NotImplementedError} from the DMRG-side driver, matching the
same scope limitation \texttt{nhdmrg()} itself doesn't have but this
feature's C++ side does --- no \texttt{Chain::nhkpm\_moments} was
ported to \texttt{mpscpp2}). The pure-Python backend's own
\texttt{Chain.nhkpm\_moments} (\texttt{pyitensor/chain.py}) is a
line-for-line transcription of \texttt{mpscpp3/chain\_session.h}'s
version, built on the same MPO/MPS algebra
(\texttt{applyMPO}/\texttt{sum}/\texttt{inner}) \texttt{general\_kpm}
already used there; no new pyitensor primitives were needed. All three
backends agree to machine precision
(\texttt{examples/non\_hermitian/nhkpm\_v3\_VS\_ED} for ED vs v3,
\texttt{examples/non\_hermitian/nhkpm\_python\_VS\_v3\_timing} for v3
vs \texttt{"python"} --- the latter on a non-uniform-hopping,
non-uniform imaginary-onsite-energy interacting fermionic chain, chosen
to avoid the Re-degenerate-ground-state pitfall a uniform staggered
pattern can trigger; \texttt{"python"} measured $\sim$1.8--2x slower
than v3 for this workload, consistent with NH-KPM's per-frequency
moment recursion being far more matvec-heavy than the Hermitian KPM
path, which amortizes its moments over the whole spectrum instead of
recomputing per frequency). Validated directly against the live
upstream Julia implementation, not just by hand: the reference's own
shipped \texttt{.OUT} output files did not reproduce under the
parameters recorded in the currently-checked-in example scripts
(plausibly stale, given the upstream repo ships multiple undated
\texttt{\_backup}/\texttt{\_old} source variants of the same
algorithm), so instead the exact
\texttt{get\_vn\_NH}/\texttt{get\_mu\_n\_NH}/\texttt{get\_spec\_kpm\_NH}/
\texttt{dos\_kpm\_NH} function bodies were run directly (Julia was
available; only the \texttt{mode="matrix"} branches were needed, which
have no ITensors.jl/Arpack.jl dependency) on the same single-particle
Hatano-Nelson-type chain used in the reference's own
\texttt{Hatano.jl} example (asymmetric-hopping non-Hermitian SSH chain,
$N=20$, all-real spectrum). Comparing the resulting tDOS$(E)$ ($E$
scanned along the real axis at fixed small imaginary broadening)
against this dmrgpy port on the identical matrix and identical
$E_{\max}$/moment count/kernel/energy grid: all 22 detected peaks land
at bit-for-bit identical energies and heights (max relative difference
$\sim3\times10^{-15}$, i.e.\ floating-point round-off, not an
algorithmic discrepancy), and every peak sits at the real part of a
genuine eigenvalue of the underlying matrix, to within the energy
grid's resolution. (Earlier, weaker checks are kept for context: the
coupled recursion also reproduces an independently-derived scalar
recursion exactly for a $1\times1$ test case, and the reconstructed
density correctly sharpens with more moments.)

\subsection{STM/Kondo tunneling spectra (\texttt{kondospectrumtk/})}

\texttt{Spin\_Chain.get\_kondo\_spectrum} (physics documented in the user
guide's ``STM/Kondo tunneling spectra'' section) is a different kind of
ED calculation from the dynamical-correlator machinery above: it needs the
\emph{full} spectrum (every eigenstate, as a possible virtual
intermediate state in the third-order perturbation theory), not just the
ground state or a resolvent at a target frequency. It reuses the ED
backend's existing \texttt{EDchain.get\_diagonalized\_hamiltonian} (dense
\texttt{eigh}, already used elsewhere for small systems) rather than
adding a new diagonalization path, and is deliberately independent of a
chain's own \texttt{itensor\_version}/mode setting -- DMRG cannot supply
the full spectrum this method needs regardless of which backend a chain
would otherwise use, so \texttt{KondoSpectrum}
(\texttt{kondospectrumtk/edkondo.py}) always goes through
\texttt{Spin\_Chain.get\_ED\_obj()} directly.

Two of the source paper's own closed-form equations for supporting
numerical functions (the temperature-broadened step and a
temperature-broadened logarithmic Kondo function) do not reproduce the
behavior the paper itself describes and plots for them;
\texttt{kondospectrumtk/stepfunctions.py} re-derives both from the
paper's own unambiguous defining integrals instead, each independently
verified against digitized values from the paper's own figures. See
that module's docstring and the user guide's ``STM/Kondo tunneling
spectra'' section for the full derivation notes and the
feature's other scope limitations (single tip-coupled site, third-order
terms are single-direction only, the potential-scattering interference
term's general-spin form is an extrapolation from the paper's own worked
example).

\textbf{\texttt{mode="DMRG"}} (\texttt{T=0} only) is a second,
independent route through this same feature that \emph{does} stay within
\texttt{itensor\_version=3}, never diagonalizing beyond the ground state
-- viable because at \texttt{T=0} both terms simplify enough to avoid
needing the full spectrum \texttt{mode="ED"} requires:

\begin{itemize}
\item The second-order term (\texttt{secondorder\_dc.py}) becomes a
  \texttt{Theta0}-weighted cumulative integral of the ordinary
  \texttt{T=0} dynamical structure factor, so it reuses
  \texttt{get\_dynamical\_correlator} (\texttt{submode="KPM"}/\texttt{"CVM"})
  as-is.
\item The third-order Kondo term (\texttt{twotime.py},
  \texttt{dmrgtwotime.py}) is a genuinely new construction: a Heisenberg
  three-point function
  \texttt{G(t2,tau)=<GS|Sl(t2+tau)Sk(t2)Sj(0)|GS>} built via real TDVP
  time evolution (a "checkpoint-and-branch" pattern: evolve
  \texttt{Sj|GS>} forward and backward through \texttt{t2}, apply
  \texttt{Sk} at each checkpoint, evolve each branch further through
  \texttt{tau}, overlap with a fixed \texttt{Sl|GS>} reference at every
  step -- reusing the single-step \texttt{tdvp\_step} primitive
  \texttt{tdz.py} already drives manually the same way), then converted
  to a frequency-domain quantity via two closed-form time-domain kernels
  (\texttt{Theta0}'s Cauchy-principal-value kernel, computed via an
  FFT-based Hilbert transform; \texttt{F0}'s kernel, smooth and
  closed-form) rather than by evaluating those functions pointwise on a
  discrete frequency grid -- an initial FFT-grid attempt at the same
  physics did not converge robustly (a discontinuous step function and a
  log singularity both have slowly-decaying spectral content that a
  finite discrete grid resolves poorly), which is why this construction
  exists in the first place rather than a simpler direct approach.
  \texttt{twotime.py}'s functions are backend-agnostic (they only ever
  consume a discretely-sampled \texttt{G(t2,tau)} array or batches
  thereof) -- \texttt{edtwotimeref.py} supplies the same interface via
  exact ED eigenbasis time evolution, used both as the development-time
  reference this construction was validated against and as a fast, exact
  ground truth for \texttt{dmrgtwotime.py}'s own test.
\end{itemize}

\texttt{dmrgtwotime.py} was written against this codebase's existing,
verified DMRG API and validated once a compiled ITensor v3 backend
became available: \texttt{G(t2,tau)} matches the ED reference to
$\sim10^{-9}$--$10^{-10}$ pointwise, and the swept third-order Kondo term
matches a grid-consistent ED reference to $\sim10^{-10}$
(\texttt{tests/test\_kondo\_spectrum\_dmrgtwotime.py}, skipped
automatically when no compiled \texttt{itensor\_version=3} backend is
available). Getting there surfaced three real bugs, none of which showed
up in the ED-only testing this module was originally written against:

\begin{itemize}
\item \texttt{tdvp\_step} silently renormalizes its output to unit norm
  on \emph{every} call. \texttt{Sj|GS>} (the state each two-time
  trajectory starts from) is generally not unit-norm (\texttt{Sj} isn't
  norm-preserving), so letting that renormalization stand discarded the
  true amplitude at every step -- confirmed directly, it produced
  overlaps exactly a factor of 2 too large end to end for a state of
  norm 0.5. Fixed by evolving a normalized copy and rescaling every
  returned checkpoint by the true original norm.
\item The forward/backward time-stepping loop shared one running
  \texttt{times} list, so the "backward" branch kept counting down from
  the forward branch's endpoint instead of from \texttt{t=0} -- it never
  actually reached negative times at all. Fixed by walking forward and
  backward from \texttt{t=0} independently and merging afterward.
\item \texttt{twotime.py}'s \texttt{t2}-integral used \texttt{np.trapz}
  per chunk, which is exactly 0 for a single-point chunk (no width to
  integrate over) -- and DMRG chunks are always single-point (each
  \texttt{t2} checkpoint is its own real trajectory), so this silently
  zeroed the entire third-order Kondo term end to end. Fixed with a
  uniform Riemann-sum accumulation (trapezoidal endpoint correction only
  at the true global first/last \texttt{t2} points) that is well defined
  for any chunk size.
\end{itemize}

A 1-site test chain also hit an unrelated internal ITensor v3 error
(building the Hamiltonian MPO, distinct from the two-site-\texttt{dmrg()}
short-chain crash \texttt{mode.py}'s ED fallback already guards against)
-- chains need at least 3 sites for this feature. The second-order term
(\texttt{submode="KPM"}) was spot-checked too, agreeing to within a few
tens of percent at thresholds, consistent with the expected
delta-broadening/moment-truncation error on top of what the ED path
already has.

The potential-interference term (\texttt{U!=0}, part of
\texttt{order=3}) is also supported for \texttt{mode="DMRG"}, via
\texttt{potentialdc.py}: its own $T=0$ limit collapses the excited-state
sum to a convolution of the \emph{same} $T=0$ dynamical structure factor
against the $F_0$ kernel instead of $\Theta_0$'s cumulative-sum
weighting ($\Theta_0(\mathrm{eV}-x)$ is a step, integrated cumulatively;
$F_0$ is not, so this piece genuinely convolves rather than cumulatively
sums) -- so it reuses \texttt{get\_dynamical\_correlator} exactly like
the second-order term above, needing no excited-state enumeration
either, and no new DMRG-side primitive. Verified (\texttt{mode="ED"},
\texttt{submode="ED"}) against the excited-state-sum
\texttt{conductance.third\_order\_potential\_dIdV} to $\sim0.07\%$
relative error (\texttt{tests/test\_kondo\_spectrum\allowbreak\_potentialdc.py}); a
real \texttt{mode="DMRG"}, \texttt{submode="KPM"} run through the full
public API was tried directly but found impractically expensive for a
routine test -- a single KPM dynamical-correlator call took $\sim40$s on
the development machine even with a trivially small number of moments,
confirmed directly -- so the
\texttt{Spin\_Chain.\_get\_kondo\_spectrum\_dmrg} wiring that combines
this term with the other two is instead checked with the three
underlying (expensive) calls monkeypatched out
(\texttt{tests/test\_kondo\_spectrum.py::\allowbreak test\_get\_kondo\_spectrum\_dmrg\allowbreak\_combines\_terms\_correctly}).

\subsection{TDZ / complex-time-evolution dynamical correlator}

\texttt{tdz.py} implements \texttt{submode="TDZ"} (Cao, Lu, Stoudenmire
\& Parcollet, ``Dynamical correlation functions from complex time
evolution,'' arXiv:2311.10909): instead of evolving along the real-time
axis (\texttt{submode="TD"}, \texttt{timedependent.py}), it evolves
along a complex-time contour $z(t,\alpha_0)$, whose negative imaginary
part damps high-energy content as the simulation proceeds, keeping the
MPS bond dimension needed for a given accuracy much smaller than under
real-time evolution alone. The true real-time correlator is then
recovered order by order via a perturbative Taylor expansion in
$\alpha_0$ around the simulated contour
(\texttt{\_reconstruct\_real\_axis}, hardcoding the paper's own explicit
$n\le4$ Appendix~B formulas). Each order needs only a fixed set of
precomputed overlap targets ($\{H^n(B|\mathrm{GS}\rangle)\}$, built once
via repeated \texttt{toMPO}/\texttt{StaticOperator} application,
independent of $t$) plus pure scalar contour integrals -- no new
tensor-network machinery beyond a single new per-step propagator
primitive:

\begin{itemize}
\item \textbf{\texttt{itensor\_version} 3 or \texttt{"python"},
\texttt{tevol\_method="TDVP"}} (the paper's own setup): a single
two-site-TDVP step with a \emph{complex} time argument --
\texttt{pyitensor/tdvp.py}'s Krylov-exponentiation core
(\texttt{\_lanczos\_expm\_multiply}) was already fully generic to
complex coefficients with no changes needed;
\texttt{mpscpp3/chain\_session.h}'s \texttt{Chain::tdvp\_step} (moved
from a private helper to a public method and widened from
\texttt{double dt} to \texttt{Cplx dt}) simply forwards to the vendored
\texttt{TDVP/tdvp.h}, which documents its own time argument as natively
``real, imaginary, or complex.''
\item \textbf{\texttt{itensor\_version=2}} (mpscpp2 has no TDVP at
all): a single MPO-Taylor step, \texttt{Chain::evolve\_taylor\_step} in
both \texttt{mpscpp2/chain\_session.h} and
\texttt{mpscpp3/chain\_session.h} (the latter as a cross-check /
non-TDVP alternative), built on the existing \texttt{evoloperator()}
Taylor expansion of $\exp(zH)$ (also mpscpp2-only, now widened from a
real \texttt{dt} to a complex \texttt{z} throughout, including
\texttt{pyitensor/chain.py}'s own \texttt{\_evoloperator}) -- the
pre-existing deliberately-reproduced ``$z^3/6$ multiplies H2 not H3''
quirk (see CLAUDE.md) is unaffected by this widening.
\end{itemize}

Current scope: only the ``greater'' branch of the correlator is
computed (the same simplification \texttt{submode="TD"} already makes),
fed into the same windowing/FFT tail as \texttt{"TD"} (factored out into
\texttt{timedependent.\_fourier\_transform\_correlator} so both
submodes share it).

\section{Backend performance: v3 vs the pure-Python backend}
\label{sec:perf}

The pure-Python backend (\texttt{itensor\_version="python"}) trades raw
speed for zero build dependencies. The tables below were originally
measured with \texttt{numba} explicitly opted in
(\texttt{pyitensor.kernels.USE\_NUMBA = True}); the \emph{default} path
(\texttt{USE\_NUMBA = False}) has since been rewritten to precompute its
own contraction plan once per bond/Lanczos-or-Krylov loop rather than
re-deriving it on every matvec call (see \texttt{pyitensor/kernels.py}'s
module docstring, ``plain path'' section), so it now performs comparably
to these numba-opted-in numbers for ground-state DMRG (spot-checked
directly: within $\sim$10\% of the table below at $n=8$--$20$) without
any opt-in at all. \texttt{USE\_NUMBA} is not just unnecessary now but
frequently counter-productive: re-measured directly, an independent-chain
workload (each chain starting from its own random MPS, e.g.\ a parameter
sweep) never reaches a stable steady state with numba on, because each
chain's own random start produces a different bond-growth trajectory and
hence different contraction shapes to compile every time --- see
\texttt{kernels.py}'s own docstring for the full, current measurements
and guidance (numba is now only ever worth opting into for a single
long-running TDVP/METTS-like session, and even there the win is modest,
$\sim$10\%, not a large one). Absolute times below will vary by machine
and load, but the qualitative trends should hold.

\subsection{Ground state energy (Heisenberg spin-1/2 chain)}

\begin{center}
\begin{tabular}{rrrr}
\toprule
$n$ & v3 (s) & python (s) & ratio (python / v3) \\
\midrule
8  & 0.10 & 0.16  & 1.6$\times$ \\
12 & 0.38 & 0.51  & 1.3$\times$ \\
16 & 0.66 & 1.12  & 1.7$\times$ \\
20 & 1.09 & 2.44  & 2.2$\times$ \\
24 & 1.58 & 4.06  & 2.6$\times$ \\
28 & 1.84 & 6.93  & 3.8$\times$ \\
30 & 2.57 & 10.84 & 4.2$\times$ \\
32 & 2.80 & 14.20 & 5.1$\times$ \\
\bottomrule
\end{tabular}
\end{center}

The ratio grows with system size: v3 benefits from real ITensor's
QN-block-sparse tensors, which \texttt{pyitensor}'s dense NumPy tensors
don't have even with numba JIT acceleration on the hot matvec loop. For
quick, small-system exploratory work the pure-Python backend is quite
competitive (1.3--2$\times$); for production-size chains ($n \gtrsim 30$)
expect it to be several times slower.

\subsection{KPM dynamical correlator ($\langle S^z_0(t) S^z_0(0)\rangle$, Heisenberg chain)}

\begin{center}
\begin{tabular}{rrrr}
\toprule
$n$ & v3 (s) & python (s) & ratio \\
\midrule
6  & 0.19 & 0.35 & 1.8$\times$ \\
8  & 0.39 & 1.58 & 4.0$\times$ \\
10 & 1.11 & 4.54 & 4.1$\times$ \\
12 & 3.63 & 8.17 & 2.2$\times$ \\
\bottomrule
\end{tabular}
\end{center}

Dynamical correlators are generally worse for the Python backend than
plain ground-state DMRG (2--4$\times$ vs 1.3--2.6$\times$ at comparable
$n$): KPM performs many moment-recursion sweeps at a fixed, comparatively
large bond dimension (\texttt{kpmmaxm=50} by default, vs
\texttt{maxm=30} for ground states), so the missing block-sparsity cost
is paid on every sweep rather than amortized. Systems around $n\approx
30$ were not benchmarked to completion here --- extrapolating from this
table, expect a KPM correlator run at that size on the pure-Python
backend to take on the order of minutes or more.

\subsection{Other calculation types (single-run timings, $n=6$--$8$)}

\begin{center}
\begin{tabular}{lrrr}
\toprule
Calculation & v3 & python & ratio \\
\midrule
Excited-state gap (TFIM, $n=8$) & 0.41s & 2.35s & 5.7$\times$ \\
Hubbard ground state (3 sites) & 0.04s & 0.11s & 3.1$\times$ \\
\bottomrule
\end{tabular}
\end{center}

TDVP-based workloads (real-time evolution, and METTS finite-temperature
sampling in particular) are a materially different, larger gap than the
ground-state/KPM numbers above, because they call the same effective-
Hamiltonian matvec thousands of times more per run (many time steps
$\times$ many Lanczos/Krylov iterations $\times$ many METTS samples, vs.\
a handful of DMRG sweeps) --- see
\texttt{examples/finite\_temperature/backend\_timing\_metts\_dynamical\_correlator/main.py},
which measures this directly (dynamical METTS, $n=10$: v3 $\sim$9s vs.\
python $\sim$51--59s for \texttt{nsamples}=10--40, a 5--6$\times$ gap
after the matvec-planning and Lanczos-bookkeeping fixes referenced above
--- it was $\sim$13$\times$ before them).

\subsection{Practical guidance}

\begin{itemize}
  \item Use \texttt{itensor\_version="python"} for quick prototyping,
    CI/environments without a C++ toolchain, or small systems ($n
    \lesssim 20$) where the 1.3--2$\times$ overhead doesn't matter.
  \item \texttt{USE\_NUMBA}/\texttt{USE\_JAX} (both off by default, see
    the note above) are no longer needed to get the numbers above ---
    the default path already precomputes its own contraction plan.
    Opting into \texttt{USE\_NUMBA} is only ever worth it for a single
    long-running TDVP/METTS-like session (same matvec shapes recurring
    many times within one run), and even there the win is modest
    ($\sim$10\%); it's actively counter-productive for independent-chain
    workloads like parameter sweeps (see \texttt{kernels.py}'s own
    docstring for the measurements behind both claims).
  \item For production-size chains, dynamical correlators, TDVP/METTS
    workloads, or anything performance-sensitive, compile the C++
    extension (\texttt{python install.py}) and use
    \texttt{itensor\_version=2} or \texttt{3}.
  \item \texttt{examples/groundstate/backend\_timing\_gs\_energy/main.py} reproduces
    the $n=8..20$ rows of the ground-state table directly against the
    current codebase and current machine ($n=24..32$ aren't included, to
    keep the example fast to run); \texttt{examples/finite\_temperature/backend\_timing\_metts\_dynamical\_correlator/main.py}
    does the same for the TDVP/METTS gap discussed above --- re-run both
    rather than trusting these numbers verbatim on a different setup.
\end{itemize}

\section{Directory reference}

\begin{center}
\begin{longtable}{@{}p{0.42\textwidth}p{0.5\textwidth}@{}}
\toprule
Path & Contents \\
\midrule
\endhead
\texttt{src/dmrgpy/} & the Python library \\
\texttt{src/dmrgpy/manybodychain.py} & \texttt{Many\_Body\_Chain}, the shared base class \\
\texttt{src/dmrgpy/infinitechain.py}, \texttt{pyitensor/idmrg.py}, \texttt{mpscpp3/chain\_session.h} & \texttt{Infinite\_Many\_Body\_Chain} and infinite DMRG (iDMRG), pyitensor + ITensor v3 (energy density only on v3) \\
\texttt{src/dmrgpy/multioperator.py}, \texttt{multioperatortk/} & backend-agnostic operator representation \\
\texttt{src/dmrgpy/mode.py}, \texttt{cppext.py} & DMRG/ED and backend dispatch \\
\texttt{src/dmrgpy/mpscpp2/}, \texttt{mpscpp3/} & vendored ITensor C++ (v2, v3) + pybind11 bindings \\
\texttt{src/dmrgpy/pyitensor/} & pure-Python ITensor-v3-subset reimplementation \\
\texttt{src/dmrgpy/mpsjulialive/}, \texttt{mpsjulia/} & Julia/ITensors.jl backend modules \\
\texttt{src/dmrgpy/edtk/}, \texttt{pyfermion/}, \texttt{pyspin/}, \texttt{pyboson/}, \texttt{pyzn/} & exact-diagonalization backend \\
\texttt{src/dmrgpy/kpmdmrg.py} & Kernel Polynomial Method dynamical correlators \\
\texttt{src/dmrgpy/nonhermitian/kpm.py} & non-Hermitian KPM dynamical correlator (NHKPM.jl port, ED + itensor\_version=3) \\
\texttt{src/dmrgpy/tdz.py} & complex-time-evolution dynamical correlator (``TDZ'', arXiv:2311.10909) \\
\texttt{examples/} & 100+ self-contained example scripts (also the regression suite) \\
\texttt{examples/backend\_comparison/}, \texttt{v2\_VS\_v3\_*} examples & backend-vs-backend correctness comparisons \\
\texttt{examples/groundstate/\allowbreak backend\_timing\_gs\_energy/} & backend-vs-backend timing comparison \\
\texttt{installtk/} & build-requirement checking and ITensor/extension compilation \\
\texttt{docs/} & this document, and its Markdown counterpart \\
\bottomrule
\end{longtable}
\end{center}

\section{Further reading}

\begin{itemize}
  \item \texttt{README.md} --- installation quick-start and further usage
    examples (static correlators, CFT central charge, bond-dimension
    convergence).
  \item \texttt{CLAUDE.md} --- a more detailed, implementation-level
    architecture reference (written for AI coding assistants working in
    this repository, but equally useful to a human maintainer diving
    into a specific subsystem).
  \item Tutorials linked from \texttt{README.md} cover many-body quantum
    magnetism, correlated fermionic systems, and tensor-network methods
    more broadly.
\end{itemize}

\end{document}
